% HiddenPower — итерация 8, финальная редакция.
% Глава: «Связи с машинным обучением».
% Опорные источники: S. Prince; N. Thuerey et al.
% Специализированные источники: A. G. Baydin et al.; M. Bronstein et al.; N. Kovachki et al.
%
% Глава связывает математическую и физическую основу проекта с обучаемыми моделями,
% суррогатами, операторным обучением и генеративным проектированием.
\chapter{Связи с машинным обучением}
\label{chap:connections-machine-learning}
Предыдущие главы построили математический и физический контур задачи проектирования. Линейная алгебра задала пространства параметров и операторов; многомерный анализ --- дифференциалы и правило цепочки; дифференциальная геометрия --- язык формы и симметрий; динамические и стохастические модели --- эволюцию состояния; уравнения в частных производных и механика сплошной среды --- вычислимый физический отклик конструкции. Теперь требуется точно определить, где в этом контуре появляется обучаемая модель и что именно она должна приближать.

Сквозным объектом остаётся параметризованная конструкция. Её описание обозначается через $\theta\in\Theta$, физическое состояние --- через $u(\theta)$, а вычисляемый критерий качества --- через $J(\theta,u)$. Машинное обучение не заменяет эти объекты общим словом ``данные''. Оно вводит дополнительное параметризованное отображение $f_w$, где $w$ --- обучаемые параметры. Смысл $f_w$ определяется только после указания входа, выхода, функции потерь и независимой проверки.

Основной вычислительный контур главы имеет вид
\begin{equation}
 \theta
 \longmapsto
 x(\theta)
 \xmapsto{\ f_w\ }
 \widehat y
 \longmapsto
 \mathcal L(\widehat y,y)
 \longmapsto
 \nabla_w\mathcal L,
\end{equation}
где $x(\theta)$ --- выбранное представление конструкции или условий задачи, $y$ --- требуемый физический отклик, а $\widehat y$ --- предсказание. Если модель включена в проектный цикл, к этой цепочке добавляется физический решатель и производная по параметрам конструкции. Поэтому малая ошибка на обучающей выборке ещё не означает физической корректности и тем более не гарантирует пригодности модели к генеративному проектированию.

Глава не является обзором популярных архитектур. Её задача --- связать уже введённую математику с пятью конкретными вопросами: какое отображение обучается, как вычисляется его градиент, как выбирается представление геометрии, чем суррогат отличается от физического оператора и как проверяется гибридный вычислительный цикл. Такая постановка соответствует типу задач, возникающих в исследованиях AIRI по машинному обучению и генеративному проектированию, но не приписывает организации конкретных внутренних методов или проектов.

\section{Обучаемое отображение, данные и функция потерь}
\label{sec:ml-map-data-loss}
Машинное обучение начинается не с выбора архитектуры, а с указания отображения, которое требуется приблизить. В физической задаче обычно существует недоступное в замкнутом виде или дорогое отображение
\begin{equation}
 f^\star:\mathcal X\to\mathcal Y,
 \qquad
 y=f^\star(x).
\end{equation}
Звёздочка подчёркивает, что речь идёт о целевом законе, а не о конкретной реализации. Значение $y$ может быть измерено в эксперименте, вычислено численным решателем или задано аналитической моделью. Обучаемая модель
\begin{equation}
 f_w:\mathcal X\to\mathcal Y,
 \qquad
 w\in\mathcal W,
\end{equation}
должна воспроизводить нужную часть действия $f^\star$ на заранее указанной области входов. Параметры $w$ принадлежат модели и не совпадают с параметрами конструкции $\theta$. Это различие сохраняется во всей главе:
\begin{equation}
 \theta\in\Theta
 \quad\text{задаёт конструкцию},
 \qquad
 w\in\mathcal W
 \quad\text{задаёт обучаемое отображение}.
\end{equation}

\subsection{Контракт обучаемой задачи}
\begin{definition}
Контрактом обучаемой задачи назовём шестёрку
\begin{equation}
\label{eq:learning-contract-six}
 \mathfrak C
 =
 \bigl(
 \mathcal X,
 \mathcal Y,
 \mathcal D,
 \ell,
 \mathcal P_{\mathrm{use}},
 \mathcal V
 \bigr),
\end{equation}
где $\mathcal X$ и $\mathcal Y$ --- пространства входов и выходов, $\mathcal D$ --- доступные данные, $\ell$ --- функция потерь, $\mathcal P_{\mathrm{use}}$ --- распределение или область предполагаемого использования, а $\mathcal V$ --- независимый протокол проверки.
\end{definition}

Запись только архитектуры $f_w$ не определяет задачу. Одна и та же сеть может приближать податливость конструкции, поле перемещений или поправку к численному решателю. Смысл появляется лишь после фиксации всех компонентов~\eqref{eq:learning-contract-six}.

Для параметризованной конструкции вход часто строится из её параметров и условий нагружения:
\begin{equation}
 x=x(\theta,q),
\end{equation}
где $q$ объединяет нагрузку, закрепления и другие условия задачи. Физический решатель задаёт состояние
\begin{equation}
 u=\mathcal S(\theta,q),
\end{equation}
после чего из него извлекается требуемый выход
\begin{equation}
 y=\mathcal O(\theta,q,u).
\end{equation}
Таким образом, целевое отображение имеет композиционный вид
\begin{equation}
 f^\star
 =
 \mathcal O\circ\mathcal S.
\end{equation}
Эта формула показывает, что обучается не ``физика вообще'', а конкретный наблюдаемый результат физического расчёта.

\begin{connection}[с генеративным проектированием]
Генератор будущей главы~9 выдаёт кандидата $\theta=G_\phi(z,c)$ по латентному коду $z$ и условию $c$. Модель $f_w$ может быстро оценивать этот кандидат, но её контракт должен совпадать с областью кандидатов генератора. Если $G_\phi$ создаёт формы, отсутствующие в обучающих данных $\mathcal D$, малая ошибка $f_w$ на старой выборке не характеризует качество проектного цикла.
\end{connection}

\subsection{Что может быть входом и выходом}
Простейший случай --- конечномерная регрессия. Вход содержит размеры, характеристики материала и нагрузку,
\begin{equation}
 x
 =
 (\theta_1,\ldots,\theta_p,q_1,\ldots,q_r)
 \in\R^{p+r},
\end{equation}
а выход является скаляром или коротким вектором:
\begin{equation}
 y
 =
 \bigl(
 J_{\mathrm{comp}},
 M,
 \sigma_{\max}
 \bigr)
 \in\R^3.
\end{equation}
Такой суррогат удобен для быстрого ранжирования множества кандидатов, но не восстанавливает распределение перемещений и напряжений внутри тела.

Для предсказания поля выходом является функция
\begin{equation}
 y=u(\cdot)\in\mathcal U.
\end{equation}
После дискретизации он записывается вектором $U_h\in\R^{n_h}$, размер которого зависит от сетки. Здесь недостаточно сказать, что модель выдаёт ``много чисел'': необходимо указать, каким точкам или базисным функциям соответствуют компоненты и как сравниваются результаты на разных сетках.

В обратной задаче роли меняются. Наблюдение становится входом, а неизвестные параметры --- выходом:
\begin{equation}
 x=y_{\mathrm{obs}},
 \qquad
 f_w(x)\approx\theta.
\end{equation}
Такое отображение может быть неоднозначным. Два разных проекта способны давать почти одинаковый измеряемый отклик. Тогда детерминированная модель, обученная по среднеквадратичной ошибке, часто возвращает среднее между допустимыми решениями, которое само может оказаться физически бессмысленным. Здесь достаточно зафиксировать неоднозначность контракта; вероятностные и генеративные представления множества решений рассматриваются в главе~9.

\subsection{Как возникает выборка}
Пусть входы рассматриваются как реализации случайной величины $X$ с распределением $P_X$, а выход удовлетворяет
\begin{equation}
 Y=f^\star(X)+\varepsilon,
\end{equation}
где $\varepsilon$ описывает шум измерений, ошибку дискретизации или неучтённые факторы. Набор данных имеет вид
\begin{equation}
 \mathcal D_N
 =
 \{(x_i,y_i)\}_{i=1}^{N}.
\end{equation}
В задачах механики пара $(x_i,y_i)$ часто создаётся расчётом:
\begin{equation}
 x_i=x(\theta_i,q_i),
 \qquad
 u_i=\mathcal S_h(\theta_i,q_i),
 \qquad
 y_i=\mathcal O_h(\theta_i,q_i,u_i).
\end{equation}
Индекс $h$ напоминает, что данные получены дискретной моделью. Поэтому модель $f_w$ сначала учится воспроизводить $\mathcal S_h$ и $\mathcal O_h$, а не непрерывную физическую систему. Ошибка исходного решателя автоматически переходит в обучающие цели.

Область генерации данных задаётся явно:
\begin{equation}
\label{eq:parameter-domain}
 \Theta_{\mathrm{data}}
 =
 \left\{
 \theta:
 a_j\leq\theta_j\leq b_j,
 \ j=1,\ldots,p
 \right\}.
\end{equation}
К этим границам добавляются ограничения допустимости: положительная толщина, невырожденная геометрия, корректная сетка, положительно определённые материальные параметры и достаточные закрепления. Случайный вектор чисел внутри прямоугольника ещё не обязательно задаёт физически допустимую конструкцию.

Данные должны покрывать не только ``типичные'' точки, но и области, важные для решения проектной задачи. Если генеративный алгоритм будет искать лёгкие конструкции около минимальной допустимой толщины, именно эта граница требует плотной и отдельной проверки. Равномерное распределение примеров по параметрам не означает равномерного распределения по физическому отклику: небольшое изменение геометрии около концентратора может резко изменить $\sigma_{\max}$.

\subsection{Функция потерь и физический масштаб}
Для скалярной регрессии стандартная квадратичная потеря равна
\begin{equation}
\label{eq:squared-loss}
 \ell_{\mathrm{sq}}(\widehat y,y)
 =
 \frac12(\widehat y-y)^2.
\end{equation}
Она удобна для дифференцирования и сильнее штрафует большие ошибки. Однако число в~\eqref{eq:squared-loss} зависит от единиц измерения. Ошибка $1$ означает разные вещи для перемещения в метрах, напряжения в паскалях и массы в килограммах.

Для векторного выхода вводится масштаб $s_k>0$ каждой компоненты:
\begin{equation}
 \ell_{\mathrm{scaled}}(\widehat y,y)
 =
 \frac12
 \sum_{k=1}^{m}
 \alpha_k
 \left(
 \frac{\widehat y_k-y_k}{s_k}
 \right)^2,
 \qquad
 \alpha_k>0.
\end{equation}
Масштаб может быть характерным физическим значением, допустимой погрешностью или статистическим разбросом по обучающей части данных. Коэффициент $\alpha_k$ задаёт важность компоненты для задачи. Его нельзя выбирать только так, чтобы графики потерь выглядели одинаково.

Для поля на конечно-элементной сетке покомпонентное среднее
\[
 \frac1{n_h}\sum_{j=1}^{n_h}
 (\widehat U_j-U_j)^2
\]
зависит от размещения узлов и не обязано приближать физическую норму. Согласованная дискретная $L^2$-ошибка имеет вид
\begin{equation}
 \ell_{L^2,h}(\widehat U,U)
 =
 \frac12
 (\widehat U-U)^{\trans}
 M_h
 (\widehat U-U),
\end{equation}
где $M_h$ --- матрица масс. Для поля напряжений средняя ошибка также не контролирует максимум. Поэтому при проектировании добавляется отдельная проверяемая величина
\begin{equation}
 e_{\sigma,\max}
 =
 \frac{
 \left|\widehat\sigma_{\max}-\sigma_{\max}\right|
 }{
 \sigma_{\mathrm{ref}}
 }.
\end{equation}
Потеря обучения и критерий окончательной проверки могут различаться. Это нормально, если различие сформулировано заранее и не скрывает опасные режимы.

Относительная ошибка
\begin{equation}
 \ell_{\mathrm{rel}}(\widehat y,y)
 =
 \frac12
 \left(
 \frac{\widehat y-y}{|y|+\epsilon}
 \right)^2
\end{equation}
требует масштаба $\epsilon>0$. Без него значения $y$ около нуля создают искусственно огромный штраф. Выбор $\epsilon$ имеет физический смысл и должен быть связан с точностью измерения или минимально значимым откликом.

\subsection{Риск, обучение и независимая проверка}
Качество модели на предполагаемом режиме использования описывает риск
\begin{equation}
\label{eq:population-risk}
 R_{\mathrm{use}}(w)
 =
 \mathbb E_{(X,Y)\sim P_{\mathrm{use}}}
 \left[
 \ell(f_w(X),Y)
 \right].
\end{equation}
Распределение $P_{\mathrm{use}}$ обычно неизвестно полностью. По обучающей выборке вычисляется эмпирический риск
\begin{equation}
\label{eq:empirical-risk}
 \widehat R_{\mathrm{train}}(w)
 =
 \frac1{N_{\mathrm{train}}}
 \sum_{i\in I_{\mathrm{train}}}
 \ell\bigl(f_w(x_i),y_i\bigr).
\end{equation}
Минимизация~\eqref{eq:empirical-risk} подбирает параметры $w$ по обучающим данным, однако не даёт оценки популяционного риска~\eqref{eq:population-risk}.

Данные разделяются на три непересекающиеся части:
\begin{equation}
 I_{\mathrm{train}}
 \mathbin{\dot\cup}
 I_{\mathrm{val}}
 \mathbin{\dot\cup}
 I_{\mathrm{test}}
 =
 \{1,\ldots,N\}.
\end{equation}
Обучающая часть определяет $w$. Проверочная часть используется для выбора архитектуры, коэффициентов регуляризации, масштаба шага и момента остановки. Тестовая часть применяется после завершения всех таких решений. Если по тестовой ошибке многократно менять модель, тестовая выборка фактически становится ещё одной проверочной и перестаёт давать независимую оценку.

В полевых данных делить следует целые физические случаи, а не отдельные узлы. Если узлы одной и той же конечно-элементной задачи одновременно попадают в обучение и тест, модель видит почти ту же геометрию, нагрузку и решение с обеих сторон. Аналогично, последовательные моменты одной траектории сильно зависимы. Корректная единица разбиения --- независимая конструкция, независимый режим нагрузки или независимый эксперимент в зависимости от будущего применения.

Для генеративного проектирования одной случайной тестовой части недостаточно. Нужны отдельные проверки на границе допустимой области и на распределении кандидатов, которое создаёт текущий генератор:
\begin{equation}
 \theta=G_\phi(Z,c),
 \qquad
 \theta\sim P_{G_\phi}.
\end{equation}
При изменении $\phi$ меняется и $P_{G_\phi}$. Поэтому суррогат может постепенно оказаться вне собственной области надёжности, даже если в начале оптимизации распределения совпадали.

\medskip
\noindent\textbf{Переобучение, интерполяция и сдвиг распределения.}
\begin{definition}
Переобучением назовём ситуацию, в которой уменьшение ошибки на обучающей части не сопровождается уменьшением ошибки на независимых данных из предполагаемого режима использования.
\end{definition}

Разрыв обобщения можно записать как
\begin{equation}
\label{eq:generalization-gap}
 \Delta_{\mathrm{gen}}
 =
 \widehat R_{\mathrm{test}}(w)
 -
 \widehat R_{\mathrm{train}}(w).
\end{equation}
Малое значение~\eqref{eq:generalization-gap} ещё не гарантирует физической корректности: обе выборки могут быть получены одним и тем же ошибочным решателем или покрывать только лёгкую часть пространства параметров.

Интерполяцией будем называть работу внутри области, поддержанной обучающими данными, а экстраполяцией --- применение вне неё. Граница здесь не обязана совпадать с покоординатным прямоугольником~\eqref{eq:parameter-domain}. Например, в данных могут встречаться большие нагрузки и тонкие элементы по отдельности, но отсутствовать их опасная комбинация.

Сдвиг распределения возникает, когда
\begin{equation}
 P_{\mathrm{train}}\neq P_{\mathrm{use}}.
\end{equation}
В проектном цикле это типично: оптимизатор намеренно ищет редкие конструкции с меньшим значением целевой функции. Он может обнаружить не физически удачную форму, а область систематической ошибки суррогата. Поэтому лучший по предсказанию кандидат обязательно пересчитывается исходным физическим решателем, а найденная пара пополняет набор данных, если требуется продолжить поиск.

Практическая проверка начинается с предельно простой задачи: модель должна почти точно воспроизвести один или несколько примеров. Такой тест не измеряет обобщение, но обнаруживает ошибки в данных, нормировке, реализации модели и функции потерь. Только после него увеличивают выборку и оценивают независимые случаи. Включение дифференцируемого решателя или генеративной модели до прохождения этой проверки усложняет диагностику без добавления информации.

\subsection{Законченный расчёт: суррогат податливости консоли}
\label{subsec:worked-compliance-surrogate}
Рассмотрим консоль длины $L$, ширины $b$ и толщины $t$, нагруженную силой $P$ на свободном конце. В линейной модели Эйлера--Бернулли прогиб конца равен
\begin{equation}
 \delta
 =
 \frac{PL^3}{3EI},
 \qquad
 I=\frac{bt^3}{12}.
\end{equation}
Податливость, определённая как работа силы $J=P\delta$, составляет
\begin{equation}
\label{eq:cantilever-compliance}
 J
 =
 \frac{4P^2L^3}{Ebt^3}.
\end{equation}
Пусть $L$, $E$ и $b$ фиксированы. После нормировки остаётся безразмерный отклик
\begin{equation}
 y(t,p)
 =
 \frac{p^2}{t^3},
\end{equation}
где $t$ и $p$ --- нормированные толщина и сила.

Сначала выберем аффинный суррогат
\begin{equation}
\label{eq:affine-compliance-surrogate}
 \widehat y
 =
 w_0+w_t t+w_p p.
\end{equation}
Обучающая выборка содержит три расчёта:
\begin{align}
 (t,p)=(1,1)&:\quad y=1,
 \label{eq:worked-train-point-one}\\
 (t,p)=(1,2)&:\quad y=4,
 \label{eq:worked-train-point-two}\\
 (t,p)=(2,1)&:\quad y=\frac18.
\end{align}
Точное совпадение на этих точках даёт систему
\begin{align}
 w_0+w_t+w_p&=1,
 \label{eq:worked-affine-system-one}\\
 w_0+w_t+2w_p&=4,
 \label{eq:worked-affine-system-two}\\
 w_0+2w_t+w_p&=\frac18.
\end{align}
Вычитая первое равенство из второго и третьего, получаем
\begin{equation}
 w_p=3,
 \qquad
 w_t=-\frac78,
 \qquad
 w_0=-\frac98.
\end{equation}
На всей обучающей выборке ошибка равна нулю:
\begin{equation}
\label{eq:worked-zero-training-risk}
 \widehat R_{\mathrm{train}}(w)=0.
\end{equation}

Проверим не использованную при обучении комбинацию $(t,p)=(2,2)$. Истинное значение равно
\begin{equation}
 y(2,2)
 =
 \frac{2^2}{2^3}
 =
 \frac12.
\end{equation}
Аффинная модель выдаёт
\begin{align}
 \widehat y(2,2)
 &=-\frac98-\frac78\cdot2+3\cdot2
 \notag\\
 &=\frac{25}{8}
 =3.125.
\end{align}
Относительная ошибка составляет
\begin{equation}
\label{eq:worked-affine-relative-error}
 \frac{|\widehat y-y|}{|y|}
 =
 \frac{\left|25/8-1/2\right|}{1/2}
 =
 \frac{21}{4}
 =5.25.
\end{equation}
Иными словами, модель идеально запомнила три значения, но ошиблась на $525\%$ в четвёртой точке. Причина не в недостатке итераций обучения, а в несоответствии класса моделей физической зависимости: формула~\eqref{eq:affine-compliance-surrogate} не содержит ни взаимодействия $p^2$, ни зависимости $t^{-3}$.

Теперь используем физически осмысленный признак
\begin{equation}
\label{eq:physics-feature}
 \psi(t,p)=\frac{p^2}{t^3}
\end{equation}
и модель
\begin{equation}
 \widehat y=a\psi(t,p).
\end{equation}
Уже первая точка $(1,1)$ даёт $a=1$, после чего модель точно воспроизводит все четыре значения. Здесь улучшение связано не с усложнением нейросети, а с использованием структуры задачи из главы~\ref{chap:continuum-mechanics}.

Этот результат имеет строгие границы. Формула~\eqref{eq:cantilever-compliance} предполагает малые деформации, линейную упругость, постоянные $E$ и $b$, балочную кинематику и заданный способ нагружения. При изменении геометрии кронштейна, появлении отверстий или концентрации напряжений признак~\eqref{eq:physics-feature} уже не является точным законом. Тогда он может использоваться как базовая модель или входной признак, но окончательная проверка возвращается к PDE-решателю.

\begin{connection}[с AIRI и генеративным проектированием]
В исследовательском контуре генератор способен выдавать тысячи кандидатов, а полный расчёт каждого кандидата дорог. Суррогат полезен как быстрый фильтр только внутри проверенной области. Расчёт~\eqref{eq:worked-zero-training-risk}--\eqref{eq:worked-affine-relative-error} показывает, почему одной обучающей ошибки недостаточно: генератор или оптимизатор быстро находит комбинации параметров, на которых неверная модель выглядит чрезмерно оптимистичной. Поэтому быстрый прогноз, оценка области применимости и независимый физический пересчёт образуют единый алгоритм, а не три взаимозаменяемые альтернативы.
\end{connection}

\medskip
\noindent\textbf{Минимальный протокол перед обучением.}
Перед реализацией модели составляется одна строка контракта:
\begin{equation}
\label{eq:one-line-learning-contract}
 \boxed{
 x(\theta,q)
 \xmapsto{\ f_w\ }
 \widehat y
 \approx
 \mathcal O_h\bigl(\theta,q,\mathcal S_h(\theta,q)\bigr)
 }.
\end{equation}
После неё фиксируются физические единицы, область $\Theta_{\mathrm{data}}$, способ получения целей, независимая единица разбиения данных, loss и итоговые метрики. Для поля дополнительно указываются сетка и норма. Для включения в генеративный цикл указывается распределение кандидатов и правило обязательного пересчёта.

Содержательно задача готова к обучению, когда по записи~\eqref{eq:one-line-learning-contract} можно однозначно ответить на четыре вопроса: какое отображение приближается, на каких входах ему предстоит работать, какая ошибка оптимизируется и каким независимым расчётом проверяется результат. Название архитектуры до этого момента вторично.

\section{Вычислительный граф и автоматическое дифференцирование}
\label{sec:computational-graph-autodiff}
После фиксации обучаемого отображения и функции потерь требуется вычислить, как малое изменение параметров модели влияет на итоговую ошибку. Для модели с большим числом параметров полный аналитический вывод быстро становится громоздким, а конечные разности требуют отдельного запуска программы для каждого направления. Автоматическое дифференцирование решает другую задачу: оно применяет правило цепочки к фактически выполненной последовательности элементарных операций и возвращает производную этой программы в выбранной точке.

Во всей секции сохраняются обозначения раздела~\ref{sec:ml-map-data-loss}. Параметры конструкции обозначаются через $\theta$, параметры обучаемой модели --- через $w$, вход --- через $x$, предсказание --- через $\widehat y$, а скалярная функция потерь --- через $\mathcal L$. Такое разделение необходимо в гибридной задаче, где одна производная берётся по весам модели, а другая --- по геометрии или материалу конструкции.

\subsection{Вычислительный граф как запись выполненной программы}
Пусть вычисление разбито на последовательность отображений
\begin{equation}
\label{eq:computational-graph-chain}
 z_0=x,
 \qquad
 z_k=g_k(z_{k-1};w_k),
 \qquad
 k=1,\ldots,m,
 \qquad
 \mathcal L=\ell(z_m,y).
\end{equation}
Переменная $z_k$ хранит промежуточный результат, а $w_k$ --- параметры операции $g_k$. Формула~\eqref{eq:computational-graph-chain} является частным случаем ориентированного ациклического графа: одна переменная может использоваться в нескольких последующих операциях, а один параметр --- встречаться в разных узлах.

\begin{definition}
Вычислительным графом называется ориентированный граф, вершины которого соответствуют значениям промежуточных переменных, а рёбра --- непосредственным зависимостям между ними в конкретном запуске программы. Локальной производной ребра $z_i\to z_j$ называется якобиан
\begin{equation}
 J_{j\leftarrow i}
 =
 \pdv{z_j}{z_i}
\end{equation}
в значениях, полученных при прямом вычислении.
\end{definition}

Граф фиксирует именно выполненные операции. Если программа содержит условный переход, производная строится по выбранной ветви. Если цикл выполнен $T$ раз, в графе появляются $T$ последовательных применений одного правила. Поэтому автоматическое дифференцирование работает не с абстрактным названием алгоритма, а с конкретной трассой вычисления.

Для композиции двух отображений $z_1=g_1(z_0)$ и $z_2=g_2(z_1)$ правило цепочки из раздела~\ref{sec:chain-rule-sensitivity} даёт
\begin{equation}
 \pdv{z_2}{z_0}
 =
 \pdv{z_2}{z_1}
 \pdv{z_1}{z_0}.
\end{equation}
Автоматическое дифференцирование не меняет это правило. Оно только организует вычисление произведений локальных якобианов так, чтобы не строить лишние матрицы.

\medskip
\noindent\textbf{Чем автоматическое дифференцирование не является.}
Пусть $F:\R^n\to\R$ задан программой. Центральная конечная разность по координате $w_j$ имеет вид
\begin{equation}
 \pdv{F}{w_j}(w)
 \approx
 \frac{F(w+h e_j)-F(w-h e_j)}{2h}.
\end{equation}
Это приближение содержит две конкурирующие ошибки. При большом $h$ существенна ошибка усечения, а при слишком малом $h$ вычитание близких чисел усиливает округление. Для градиента по $p$ параметрам требуется порядка $2p$ дополнительных запусков программы. Поэтому конечная разность полезна как независимая проверка нескольких компонент, но не как основной способ обучения модели с большим $p$.

Символьное дифференцирование строит новое аналитическое выражение. Оно удобно для коротких формул, но при длинной композиции может многократно повторять одинаковые подвыражения. Автоматическое дифференцирование вместо этого сохраняет промежуточные численные значения и распространяет по ним производные. Результат является числом или произведением с якобианом в заданной точке, а не обязательной замкнутой формулой.

\begin{definition}
Автоматическим дифференцированием называется вычисление производных программы путём разложения её на элементарные операции с известными локальными производными и последующего применения правила цепочки к значениям, полученным при исполнении программы.
\end{definition}

При гладких элементарных операциях AD вычисляет производную с точностью машинной арифметики. Это не означает, что любая программа гладкая. Операции сравнения, округления, изменения дискретной топологии и остановки итераций могут создавать разрывы или кусочно-гладкие зависимости. Тогда вычислительная система возвращает производную выбранной ветви либо требует явно заданного правила.

\subsection{Прямой режим и произведение якобиана на вектор}
Рассмотрим отображение
\begin{equation}
 F:\R^n\to\R^m,
 \qquad
 y=F(x).
\end{equation}
Вместо полного якобиана $J_F(x)\in\R^{m\times n}$ часто нужна производная в одном направлении $v\in\R^n$:
\begin{equation}
 \dot y
 =
 J_F(x)v.
\end{equation}
Точка над переменной обозначает касательное приращение, а не производную по физическому времени. Входу сопоставляется пара $(x,\dot x)$, где $\dot x=v$. Каждая элементарная операция одновременно вычисляет значение и его направленную производную.

Для скалярных операций
\begin{equation}
 c=a+b,
 \qquad
 d=ab,
 \qquad
 r=\varphi(d)
\end{equation}
прямой режим распространяет
\begin{equation}
 \dot c=\dot a+\dot b,
 \qquad
 \dot d=\dot a\,b+a\,\dot b,
 \qquad
 \dot r=\varphi'(d)\dot d.
\end{equation}
После одного прохода получается $J_F(x)v$ без явного построения $J_F(x)$.

Если требуется весь якобиан, можно последовательно подавать базисные направления $e_1,\ldots,e_n$. Стоимость тогда растёт пропорционально числу входных координат $n$. Поэтому прямой режим особенно удобен, когда входных направлений мало, а выходов много. Типичный пример --- чувствительность большого поля перемещений к одному материальному параметру.

\begin{connection}[с геометрией конструкции]
Пусть координаты узлов сетки зависят от одного проектного параметра $\theta_j$, а решатель или суррогат возвращает поле $U_h(\theta)$. Задав $\dot\theta=e_j$, прямой режим вычисляет
\begin{equation}
 \dot U_h
 =
 \pdv{U_h}{\theta_j}.
\end{equation}
Это поле показывает, в каких узлах и направлениях изменится физическое состояние при малом изменении одного размера конструкции.
\end{connection}

\subsection{Обратный режим и произведение вектора на якобиан}
В обучении выходом вычислительного графа обычно является один скаляр $\mathcal L$. Требуется градиент по большому числу параметров. В этой ситуации выгодно начать с выхода и распространять чувствительность назад.

Для каждой промежуточной переменной $z_k$ введём ковектор
\begin{equation}
 \bar z_k
 =
 \pdv{\mathcal L}{z_k}.
\end{equation}
Если $z_{k+1}=g_{k+1}(z_k)$, то
\begin{equation}
 \bar z_k
 =
 \bar z_{k+1}
 \pdv{z_{k+1}}{z_k}.
\end{equation}
Это произведение вектора на якобиан, или VJP. Для столбцовой записи градиентов эквивалентная формула имеет вид
\begin{equation}
 \nabla_{z_k}\mathcal L
 =
 \left(\pdv{z_{k+1}}{z_k}\right)^{\trans}
 \nabla_{z_{k+1}}\mathcal L.
\end{equation}

Если переменная $z_k$ используется в нескольких последующих узлах, вклады складываются:
\begin{equation}
 \bar z_k
 =
 \sum_{j:\,z_k\to z_j}
 \bar z_j
 \pdv{z_j}{z_k}.
\end{equation}
Это обычное дифференцирование суммы всех путей, по которым $z_k$ влияет на loss. Именно поэтому повторное использование параметра не требует отдельного ручного вывода: граф накапливает все его вклады.

Один обратный проход возвращает производные скалярной функции по всем сохранённым входам и параметрам. Его арифметическая стоимость обычно сопоставима с несколькими прямыми вычислениями, но требуется доступ к промежуточным значениям. Следовательно, обратный режим меняет вычислительное время на память.

\begin{definition}
Backpropagation называется применение обратного режима автоматического дифференцирования к вычислительному графу модели и скалярной функции потерь. Это не отдельное правило дифференцирования, а специальная организация правила цепочки.
\end{definition}

\subsection{Производные по весам и по входу}
Рассмотрим один аффинный слой
\begin{equation}
 z=Wx+b,
 \qquad
 z\in\R^r,
 \quad
 x\in\R^n.
\end{equation}
Пусть из последующей части графа пришёл столбец $\bar z=\nabla_z\mathcal L$. Локальные VJP дают
\begin{equation}
\label{eq:affine-layer-backward}
 \nabla_x\mathcal L=W^{\trans}\bar z,
 \qquad
 \nabla_W\mathcal L=\bar z x^{\trans},
 \qquad
 \nabla_b\mathcal L=\bar z.
\end{equation}
Первая производная нужна для продолжения обратного прохода, а две остальные --- для обновления параметров слоя.

Формула~\eqref{eq:affine-layer-backward} показывает связь с линейной алгеброй главы~\ref{chap:linear-algebra}. Обратный проход по входу использует транспонированный линейный оператор, а производная по матрице имеет ранг не выше единицы для одного объекта. При обработке мини-пакета такие внешние произведения суммируются или усредняются.

Производная по входу имеет самостоятельный смысл. Если $x=x(\theta,q)$ кодирует конструкцию, то
\begin{equation}
 \nabla_\theta\mathcal L
 =
 \left(\pdv{x}{\theta}\right)^{\trans}
 \nabla_x\mathcal L.
\end{equation}
В обучении обычно изменяется $w$, а в генеративном проектировании после обучения может изменяться $\theta$ или параметр генератора $\phi$. Один и тот же вычислительный граф допускает разные конечные переменные дифференцирования, но их физический смысл нельзя смешивать.

\subsection{Законченный расчёт: прямой и обратный проход}
\label{subsec:worked-forward-backward}
Рассмотрим малую модель скалярного отклика конструкции:
\begin{equation}
\label{eq:worked-model}
 z_1=w_1x+b_1,
 \qquad
 z_2=z_1^2,
 \qquad
 \widehat y=w_2z_2+b_2,
 \qquad
 \mathcal L=\frac12(\widehat y-y)^2.
\end{equation}
Здесь $x$ может быть нормированным геометрическим признаком, а $\widehat y$ --- нормированной податливостью. Возьмём
\begin{equation}
 x=2,
 \qquad
 y=3,
 \qquad
 w_1=\frac12,
 \qquad
 b_1=1,
 \qquad
 w_2=2,
 \qquad
 b_2=-1.
\end{equation}

Прямой проход даёт
\begin{equation}
 z_1=\frac12\cdot2+1=2,
 \qquad
 z_2=2^2=4,
 \qquad
 \widehat y=2\cdot4-1=7,
 \qquad
 \mathcal L=\frac12(7-3)^2=8.
\end{equation}

Обратный проход начинается с
\begin{equation}
 \bar{\mathcal L}=1,
 \qquad
 \bar{\widehat y}
 =
 \pdv{\mathcal L}{\widehat y}
 =
 \widehat y-y
 =4.
\end{equation}
Из аффинной операции $\widehat y=w_2z_2+b_2$ получаем
\begin{equation}
 \pdv{\mathcal L}{w_2}
 =
 \bar{\widehat y}\,z_2
 =4\cdot4=16,
 \qquad
 \pdv{\mathcal L}{b_2}=4,
 \qquad
 \bar z_2
 =
 \bar{\widehat y}\,w_2
 =8.
\end{equation}
Далее $z_2=z_1^2$, поэтому
\begin{equation}
 \bar z_1
 =
 \bar z_2\,2z_1
 =8\cdot4
 =32.
\end{equation}
Наконец, из $z_1=w_1x+b_1$ следуют
\begin{equation}
 \pdv{\mathcal L}{w_1}
 =
 \bar z_1 x
 =32\cdot2=64,
 \qquad
 \pdv{\mathcal L}{b_1}=32,
 \qquad
 \pdv{\mathcal L}{x}
 =
 \bar z_1w_1
 =16.
\end{equation}
Итоговый градиент по параметрам модели равен
\begin{equation}
 \nabla_w\mathcal L
 =
 \left(
 64,
 32,
 16,
 4
 \right)^{\trans},
 \qquad
 w=(w_1,b_1,w_2,b_2).
\end{equation}

Расчёт показывает, что обратный проход не ``угадывает'' вклад параметра. Он последовательно умножает входящую чувствительность на локальную производную каждого узла. Значение $\partial\mathcal L/\partial x=16$ одновременно сообщает, как изменится loss при малом изменении входного геометрического признака.

\subsection{Проверка градиента конечной разностью}
Автоматическое дифференцирование устраняет ручное раскрытие длинной композиции, но не гарантирует правильность самой программы. Ошибка может находиться в реализации модели, форме loss, преобразовании единиц или пользовательском правиле производной. Поэтому перед крупным запуском несколько компонент градиента проверяются независимо.

Для параметра $w_j$ используем относительное расхождение
\begin{equation}
 \delta_j(h)
 =
 \frac{
 \left|
 g_j^{\mathrm{AD}}
 -
 \dfrac{\mathcal L(w+h e_j)-\mathcal L(w-h e_j)}{2h}
 \right|
 }{
 1+|g_j^{\mathrm{AD}}|
 }.
\end{equation}
Единица в знаменателе не даёт делить на почти нулевую производную.

В расчёте~\eqref{eq:worked-model} при фиксированных остальных величинах
\begin{equation}
 \mathcal L(w_2)
 =
 \frac12(4w_2-4)^2
 =8(w_2-1)^2.
\end{equation}
Поэтому при $w_2=2$
\begin{equation}
 \frac{\mathcal L(2+h)-\mathcal L(2-h)}{2h}
 =16
 =
 \pdv{\mathcal L}{w_2}.
\end{equation}
Для этой квадратичной зависимости центральная разность совпадает с точной производной в точной арифметике. В общей программе следует проверить несколько значений $h$: сначала ошибка убывает примерно как $h^2$, затем перестаёт уменьшаться из-за округления.

Проверка должна выполняться на малой задаче с фиксированными входами. Нельзя одновременно менять случайный мини-пакет, сетку или стохастический шум: тогда сравниваются разные функции. Для негладкой точки совпадение также не обязано наблюдаться.

\subsection{Дифференцируемый физический оператор}
Физический решатель удобно включать в граф не как тысячи низкоуровневых операций, а как содержательный оператор
\begin{equation}
 U_{k+1}=\mathcal P_k(U_k,\theta,q).
\end{equation}
Для обратного прохода оператор должен уметь вычислять VJP
\begin{equation}
 \left(\pdv{\mathcal P_k}{U_k}\right)^{\trans}\lambda_{k+1},
 \qquad
 \left(\pdv{\mathcal P_k}{\theta}\right)^{\trans}\lambda_{k+1},
\end{equation}
а не обязательно хранить полный якобиан. Это тот же принцип, который использовался для сопряжённой чувствительности ODE в разделе~\ref{sec:sensitivity-and-adjoint-gradient} и PDE в разделе~\ref{sec:differentiable-pde-solver}.

Для неявного линейного решателя
\begin{equation}
 K(\theta)U(\theta)=F(\theta)
\end{equation}
необязательно дифференцировать все итерации численного алгоритма. Дифференцирование самого уравнения даёт
\begin{equation}
 K\,\dot U
 =
 \dot F-\dot K\,U.
\end{equation}
Если loss зависит от $U$, сопряжённая система
\begin{equation}
 K^{\trans}\lambda
 =
 \nabla_U\mathcal L
\end{equation}
позволяет получить производную по большому числу проектных параметров без отдельного решения для каждого $\dot U$. Автоматическое дифференцирование и неявное дифференцирование здесь не конкурируют: AD собирает производные внешнего графа, а физический оператор предоставляет устойчивое локальное правило VJP.

\begin{connection}[с AIRI и генеративным проектированием]
В полном контуре
\begin{equation}
 z
 \xmapsto{\ G_\phi\ }
 \theta
 \xmapsto{\ \mathcal S\ }
 U
 \longmapsto
 J(\theta,U)
\end{equation}
обратный режим вычисляет $\nabla_\phi J$ как последовательность VJP генератора, геометрического представления и физического решателя. Содержательная задача исследования состоит не в наличии команды \texttt{backward}, а в корректности каждого локального правила: геометрия не должна терять допустимость, решатель должен возвращать физически согласованную чувствительность, а критерий --- выражать реальное требование к конструкции.
\end{connection}

\medskip
\noindent\textbf{Память, повторное вычисление и протокол использования.}
Обратный режим требует промежуточные значения прямого прохода. Если граф содержит много слоёв или временных шагов, хранение всех $z_k$ становится основным ограничением памяти. Полное освобождение памяти заставило бы повторить весь прямой расчёт при каждом локальном шаге назад. Практический компромисс --- контрольные точки: сохраняется часть состояний, а промежуточные участки при необходимости пересчитываются.

Разбиение физического решателя на содержательные операции также уменьшает память. Вместо хранения каждой внутренней арифметической команды сохраняются состояния, необходимые для локального VJP выбранного шага или неявного оператора. Такое укрупнение допустимо только вместе с проверенным правилом производной; объявление сложного блока ``дифференцируемым'' без теста градиента не является достаточным.

Перед использованием производной в обучении или проектировании фиксируется цепочка
\begin{equation}
 \boxed{
 \begin{gathered}
 \text{прямой граф}
 \to
 \text{локальные JVP/VJP}
 \to
 \text{обратный градиент}
 \\
 \text{проверка конечной разностью}
 \to
 \text{физическая проверка направления}
 \end{gathered}
 }.
\end{equation}
Последний этап отличается от численной проверки. Даже математически верный градиент может вести к неприемлемой конструкции, если loss не учитывает прочность, технологичность или область применимости модели. Поэтому производная проверяется сначала как производная программы, а затем как направление изменения реального проектного критерия.

\section{Обучение как задача оптимизации}
\label{sec:training-as-optimization}
После построения функции потерь и вычисления её градиента обучение сводится к последовательному изменению параметров модели. Эта формулировка не означает, что любой оптимизатор автоматически найдёт физически содержательную модель. Алгоритм видит только выбранный функционал, масштаб переменных и данные, поэтому качество шага определяется всей постановкой из раздела~\ref{sec:ml-map-data-loss}, а не одним названием метода.

Во всей секции $w\in\mathcal W\subseteq\R^p$ обозначает параметры обучаемой модели, $\theta$ --- параметры конструкции, $\widehat R_N(w)$ --- эмпирический риск, а $\mathcal L_B(w)$ --- loss на мини-пакете $B$. Параметры $w$ обучаются по данным; параметры $\theta$ изменяются уже внутри проектного или генеративного цикла. Эти две оптимизации могут использовать похожие формулы, но решают разные задачи и не должны смешиваться в обозначениях.

\subsection{Целевая функция и область допустимых параметров}
Базовая задача обучения имеет вид
\begin{equation}
 w_*
 \in
 \argmin_{w\in\mathcal W}
 \Phi_N(w),
 \qquad
 \Phi_N(w)
 =
 \widehat R_N(w)+\lambda\mathcal R(w),
\end{equation}
где $\mathcal R$ --- регуляризатор, а $\lambda\geq0$ задаёт его вес. Символ $\argmin$ обозначает множество минимизаторов: единственность решения заранее не предполагается.

Если $\mathcal W=\R^p$ и $\Phi_N$ дифференцируема, необходимое условие внутреннего локального минимума записывается как
\begin{equation}
 \nabla_w\Phi_N(w_*)=0.
\end{equation}
Это условие не является достаточным и не утверждает, что найден глобальный минимум. Для нелинейной модели функция $\Phi_N$ обычно невыпукла, поэтому алгоритм строит последовательность приближений, а результат оценивается по проверочной выборке и физическим тестам.

Один шаг первого порядка имеет форму
\begin{equation}
 w^{k+1}
 =
 w^k-\eta_k d_k,
\end{equation}
где $\eta_k>0$ --- длина шага, а $d_k$ --- вычисленное направление. В обычном градиентном спуске $d_k=\nabla\Phi_N(w^k)$. Если используются импульс, адаптивное масштабирование или проекция, направление уже не совпадает с текущим градиентом.

\begin{connection}[с предыдущими главами]
Локальное поведение $\Phi_N$ определяется формулой Тейлора второго порядка из раздела~\ref{sec:second-order-implicit-constrained}:
\begin{equation}
\label{eq:training-local-quadratic}
 \Phi_N(w+\delta w)
 \approx
 \Phi_N(w)
 +\nabla\Phi_N(w)^{\trans}\delta w
 +\frac12\delta w^{\trans}H(w)\delta w,
\end{equation}
где $H(w)=\Hessian{\Phi_N}(w)$. Градиент задаёт линейное уменьшение, а Гессиан определяет кривизну, допустимый масштаб шага и чувствительность к выбору координат.
\end{connection}

\subsection{Полный градиент и мини-пакет}
Пусть обучающая часть содержит $N$ пар $(x_i,y_i)$. Эмпирический риск равен
\begin{equation}
 \widehat R_N(w)
 =
 \frac1N\sum_{i=1}^{N}\ell_i(w),
 \qquad
 \ell_i(w)
 =
 \ell\bigl(f_w(x_i),y_i\bigr).
\end{equation}
Полный градиент
\begin{equation}
 g(w)
 =
 \nabla\widehat R_N(w)
 =
 \frac1N\sum_{i=1}^{N}\nabla\ell_i(w)
\end{equation}
использует всю обучающую выборку. Один такой шаг дорог, если каждый пример содержит большое поле или был получен физическим решателем.

Для мини-пакета $B_k\subset\{1,\ldots,N\}$ размера $b$ вводится
\begin{equation}
 \widehat g_k
 =
 \frac1b\sum_{i\in B_k}\nabla\ell_i(w^k).
\end{equation}
Если индексы выбираются равновероятно, то при фиксированном $w^k$
\begin{equation}
 \mathbbE\bigl[\widehat g_k\mid w^k\bigr]
 =
 g(w^k).
\end{equation}
Равенство математических ожиданий не означает, что отдельный шаг совпадает с полным градиентом. Разность
\begin{equation}
\label{eq:minibatch-noise}
 \xi_k
 =
 \widehat g_k-g(w^k)
\end{equation}
является шумом оценки. При независимом выборе примеров его типичный масштаб уменьшается с ростом $b$, но стоимость одного шага растёт.

В физических данных примеры часто неоднородны. Конструкции около ограничения прочности, редкие граничные условия и разные семейства сеток нельзя оставлять на волю случайного перемешивания. Мини-пакет должен сохранять заранее выбранную схему покрытия, иначе большинство шагов будет определяться лёгкими типичными случаями, а критическая область окажется почти незаметной.

\medskip
\noindent\textbf{Длина шага и условие уменьшения.}
Рассмотрим гладкую функцию с локально ограниченной кривизной. Для шага $w^+=w-\eta\nabla\Phi_N(w)$ квадратичная модель~\eqref{eq:training-local-quadratic} даёт
\begin{equation}
 \Phi_N(w^+)-\Phi_N(w)
 \approx
 -\eta\norm{\nabla\Phi_N(w)}_2^2
 +\frac{\eta^2}{2}
 \nabla\Phi_N(w)^{\trans}H(w)\nabla\Phi_N(w).
\end{equation}
Первое слагаемое уменьшает функцию, второе ограничивает допустимый шаг. Если $\lambda_{\max}$ --- наибольшее собственное значение симметричного положительно определённого Гессиана, то для квадратичной задачи устойчивое уменьшение требует
\begin{equation}
\label{eq:quadratic-step-stability}
 0<\eta<\frac{2}{\lambda_{\max}}.
\end{equation}
Слишком большой шаг перескакивает через минимум и может вызвать расходимость. Слишком малый шаг сохраняет устойчивость, но делает обучение практически неподвижным в направлениях малой кривизны.

Для мини-пакетного градиента длина шага взаимодействует с шумом~\eqref{eq:minibatch-noise}. Даже если среднее направление верно, большой $\eta_k$ превращает разброс градиента в сильный разброс параметров. Поэтому уменьшение training loss на отдельном пакете не является достаточной диагностикой; контролируются сглаженная обучающая кривая, проверочная ошибка и физические метрики.

\subsection{Масштабирование и обусловленность}
Пусть локальная квадратичная модель имеет Гессиан $H\succ0$ с собственными значениями
\begin{equation}
 0<\lambda_{\min}\leq\lambda_i\leq\lambda_{\max}.
\end{equation}
Число
\begin{equation}
 \kappa(H)
 =
 \frac{\lambda_{\max}}{\lambda_{\min}}
\end{equation}
характеризует различие кривизны по направлениям. При большом $\kappa(H)$ шаг ограничивается крутым направлением, а вдоль пологого направления прогресс остаётся медленным. Это та же причина нестабильности, которая в линейной алгебре проявлялась в плохо обусловленных системах и требовала масштабирования или регуляризации.

Пусть физический вход содержит длину в метрах, модуль Юнга в паскалях и безразмерную плотность материала. Без нормировки одинаковое изменение числового значения этих компонент имеет разный физический смысл. Вводится аффинное преобразование
\begin{equation}
 \widetilde x_j
 =
 \frac{x_j-\mu_j}{s_j},
 \qquad
 s_j>0.
\end{equation}
Оно не добавляет данных и не исправляет неверную модель, но меняет координаты задачи и масштаб производных. Значения $\mu_j$ и $s_j$ вычисляются только по обучающей части и затем без изменения применяются к validation и test.

Для проектных параметров масштаб выбирается с учётом допустимого диапазона и физической чувствительности. Нормировка толщины и координаты узла по одной статистической формуле не гарантирует одинаковой значимости. Если одна переменная изменяет массу на доли процента, а другая переводит конструкцию через потерю устойчивости, это различие должно остаться видимым в loss и протоколе проверки.

\medskip
\noindent\textbf{Momentum и Adam как преобразование шага.}
Momentum вводит накопленное направление
\begin{align}
 v^{k+1}
 &=
 \beta v^k+(1-\beta)\widehat g_k,
 \\
 w^{k+1}
 &=
 w^k-\eta v^{k+1},
 \qquad
 0\leq\beta<1.
\end{align}
Если последовательные градиенты согласованы, компонента направления накапливается. Если мини-пакетный шум меняет знак, усреднение частично его подавляет. Momentum не устраняет неверное масштабирование и не гарантирует монотонного уменьшения loss; он изменяет динамику итераций.

Adam дополнительно оценивает покоординатный второй момент градиента:
\begin{align}
 m^{k+1}
 &=
 \beta_1m^k+(1-\beta_1)\widehat g_k,
 \\
 q^{k+1}
 &=
 \beta_2q^k+(1-\beta_2)(\widehat g_k\odot\widehat g_k).
\end{align}
После поправки начального смещения
\begin{equation}
 \widehat m^{k+1}
 =
 \frac{m^{k+1}}{1-\beta_1^{k+1}},
 \qquad
 \widehat q^{k+1}
 =
 \frac{q^{k+1}}{1-\beta_2^{k+1}},
\end{equation}
шаг записывается как
\begin{equation}
 w^{k+1}
 =
 w^k
 -\eta
 \frac{\widehat m^{k+1}}
 {\sqrt{\widehat q^{k+1}}+\varepsilon},
\end{equation}
где операции выполняются покомпонентно. Знаменатель уменьшает шаг в координатах с устойчиво большим квадратом градиента. Это адаптивная диагональная перенормировка, а не вычисление полного обратного Гессиана.

Выбор между обычным SGD, momentum и Adam проверяется на одной и той же разбивке данных и по одинаковым физическим метрикам. Нельзя сравнивать оптимизаторы только по минимальному training loss, если один из них хуже переносится на новые конструкции или систематически нарушает граничные условия.

\subsection{Регуляризация и ограничения решают разные задачи}
Регуляризатор изменяет целевую функцию. Для квадратичного штрафа
\begin{equation}
 \Phi_N(w)
 =
 \widehat R_N(w)
 +\frac{\lambda}{2}\norm{w}_2^2
\end{equation}
градиент получает добавку
\begin{equation}
 \nabla\Phi_N(w)
 =
 \nabla\widehat R_N(w)+\lambda w.
\end{equation}
Такой член сдерживает рост параметров и связан с регуляризацией плохо обусловленных задач из главы~\ref{chap:linear-algebra}. Однако малые веса сами по себе не означают физически допустимое предсказание.

Ограничение задаёт множество разрешённых решений:
\begin{equation}
 c_j(w)\leq0,
 \qquad
 h_r(w)=0.
\end{equation}
Замена ограничения штрафом
\begin{equation}
 \Phi_{N,\rho}(w)
 =
 \Phi_N(w)
 +\rho\sum_j\bigl[c_j(w)_+\bigr]^2
 +\rho\sum_r h_r(w)^2
\end{equation}
создаёт новую безусловную задачу. При конечном $\rho$ нарушение может остаться ненулевым, а слишком большой $\rho$ ухудшает обусловленность. Поэтому равновесие, положительная толщина или точное закрепление нельзя объявлять выполненными только потому, что соответствующий штраф включён в loss.

В гибридной модели жёсткое условие лучше встраивать в параметризацию или физический оператор, когда это возможно. Мягкий штраф используется для допускаемой погрешности и всегда сопровождается численным порогом проверки.

\medskip
\noindent\textbf{Остановка по проверочной ошибке.}
Итерация $k$ выбирается не по тестовой выборке. После каждой заранее установленной группы шагов вычисляется проверочная метрика
\begin{equation}
 V_k
 =
 \frac1{|I_{\mathrm{val}}|}
 \sum_{i\in I_{\mathrm{val}}}
 \ell_i(w^k).
\end{equation}
Сохраняется состояние с наименьшим $V_k$ либо с наилучшим заранее определённым набором физических показателей. Тестовая часть применяется один раз после фиксации модели и всех гиперпараметров.

Ранняя остановка является способом выбора итерации, а не доказательством обобщающей способности. Если validation и training получены из одной узкой области параметров, обе ошибки могут быть малы, а модель всё равно провалится на новой геометрии. Поэтому к $V_k$ добавляются раздельные проверки на границе допустимой области, на других сетках и на физически критических режимах.

Для воспроизводимости фиксируются начальное состояние, разбиение данных, последовательность мини-пакетов, правило изменения $\eta_k$, критерий остановки и сохранённая итерация. Без этих данных сравнение двух обучений не позволяет понять, изменился метод или только случайная траектория.

\subsection{Законченный расчёт: масштаб переменных и скорость спуска}
\label{subsec:worked-conditioning-gradient-descent}
Рассмотрим квадратичную функцию
\begin{equation}
 \Phi(w_1,w_2)
 =
 \frac12\left(w_1^2+100w_2^2\right).
\end{equation}
Её градиент и Гессиан равны
\begin{equation}
 \nabla\Phi(w)
 =
 \begin{pmatrix}
 w_1\\100w_2
 \end{pmatrix},
 \qquad
 H
 =
 \begin{pmatrix}
 1&0\\0&100
 \end{pmatrix},
 \qquad
 \kappa(H)=100.
\end{equation}
Начнём с $w^0=(1,1)^{\trans}$. Условие~\eqref{eq:quadratic-step-stability} требует $0<\eta<0.02$. Возьмём $\eta=0.01$. Тогда
\begin{equation}
 w^1
 =
 \begin{pmatrix}
 1\\1
 \end{pmatrix}
 -0.01
 \begin{pmatrix}
 1\\100
 \end{pmatrix}
 =
 \begin{pmatrix}
 0.99\\0
 \end{pmatrix}.
\end{equation}
Крутая координата исчезла за один шаг, но первая координата далее уменьшается как
\begin{equation}
 w_1^k
 =
 0.99^k.
\end{equation}
Чтобы получить $|w_1^k|\leq0.01$, требуется
\begin{equation}
 k
 \geq
 \frac{\log 0.01}{\log 0.99}
 \approx
 458.
\end{equation}
Малый шаг выбран не из-за первой координаты, а из-за большой кривизны по $w_2$.

Теперь введём новую координату $u_2=10w_2$ и $u_1=w_1$. Тогда
\begin{equation}
 \Phi
 =
 \frac12\left(u_1^2+u_2^2\right),
 \qquad
 \nabla_u\Phi
 =
 \begin{pmatrix}
 u_1\\u_2
 \end{pmatrix},
 \qquad
 \kappa(H_u)=1.
\end{equation}
При $u^0=(1,10)^{\trans}$ и $\eta=1$ получаем
\begin{equation}
 u^1
 =
 u^0-\nabla_u\Phi(u^0)
 =
 0.
\end{equation}
В этой квадратичной задаче одна перенормировка координат устранила различие кривизны и позволила дойти до минимума за один шаг. В реальной сети точный Гессиан неизвестен и меняется по траектории, но расчёт показывает причину нормировки признаков и адаптивного масштаба шага. Они не являются косметической обработкой данных: они меняют геометрию оптимизационной задачи.

\subsection{Протокол обучения физического суррогата}
Для суррогата конструкции обучение фиксируется как единый вычислительный протокол:
\begin{equation}
 \boxed{
 \begin{gathered}
 \text{масштабы и разбиение данных}
 \to
 \text{модель и оптимизатор}
 \to
 \text{mini-batch шаги}
 \\
 \to
 \text{validation и остановка}
 \to
 \text{независимый test}
 \\
 \to
 \text{физическая верификация}
 \end{gathered}
 }.
\end{equation}
На каждом этапе сохраняются не только loss, но и физически читаемые ошибки: относительная ошибка перемещений, ошибка энергии, невязка равновесия, нарушение граничных условий и ошибка проектного критерия. Низкая средняя ошибка поля не заменяет проверку максимального напряжения или податливости.

\begin{connection}[с AIRI и генеративным проектированием]
В контуре
\begin{equation}
 \underbrace{D_{\mathrm{train}}
 \longrightarrow
 w_*}_{\text{обучение суррогата}}
 \qquad
 \text{и затем}
 \qquad
 \underbrace{z
 \xmapsto{G_\phi}
 \theta
 \longrightarrow
 f_{w_*}(\theta)
 \longrightarrow
 J}_{\text{поиск конструкции}}
\end{equation}
параметры $w_*$ после обучения фиксированы, а изменяются $\phi$ или $\theta$. Такое разделение необходимо для задач генеративного проектирования, к которым готовит HiddenPower: иначе обновление суррогата и обновление конструкции смешиваются, и улучшение внутреннего loss нельзя отличить от реального улучшения физического кандидата.
\end{connection}

Итог секции состоит не в выборе одного универсального оптимизатора. Корректное обучение требует согласовать функцию потерь, масштаб координат, способ оценки градиента, длину шага, регуляризацию, критерий остановки и независимую физическую проверку. Только после этого оптимизационная траектория становится воспроизводимой частью полного вычислительного контура.

\section{Представление геометрии, сетки и конструкции}
\label{sec:geometry-mesh-representation}
Обучаемая модель получает не физическую конструкцию как таковую, а её конечную запись. Поэтому выбор представления входит в математическую постановку наравне с функцией потерь и алгоритмом обучения. Две записи могут описывать один объект, но иметь разную размерность, нумерацию узлов или систему координат. Если модель реагирует на эти служебные различия, она учит свойства файла, а не свойства конструкции.

Обозначим множество физических конструкций через $\mathcal C$, пространство их машинных представлений через $\mathcal X$, а кодирование через
\begin{equation}
 \mathsf E:\mathcal C\to\mathcal X,
 \qquad
 c\longmapsto x=\mathsf E(c).
\end{equation}
Обратное отображение может быть неоднозначным или определяться только алгоритмом построения геометрии. Поэтому далее различаются физический объект $c$, его параметры $\theta$, дискретная геометрия $\mathcal M$ и тензор входных признаков $x$. Такое разделение не позволяет незаметно заменить конструкцию одним удобным массивом.

\subsection{Представление как часть контракта модели}
Пусть две записи $x,x'\in\mathcal X$ описывают один физический объект. Запишем это как отношение эквивалентности
\begin{equation}
 x\sim x'.
\end{equation}
Для скалярного свойства конструкции, например массы или податливости, корректный предиктор должен удовлетворять
\begin{equation}
 x\sim x'
 \quad\Longrightarrow\quad
 f_w(x)=f_w(x').
\end{equation}
Если выходом является поле на узлах, равенство заменяется согласованным преобразованием компонент. Таким образом, вопрос об инвариантности возникает до выбора архитектуры: сначала задаётся, какие изменения записи не меняют физический смысл, и только затем строится модель, уважающая это отношение.

Полный контракт представления удобно записать как
\begin{equation}
 \boxed{
 c
 \xmapsto{\ \mathsf E\ }
 x
 \xmapsto{\ f_w\ }
 \widehat y,
 \qquad
 x\sim x'
 \Rightarrow
 \widehat y' = \rho_{\mathcal Y}(x,x')\widehat y
 }.
\end{equation}
Оператор $\rho_{\mathcal Y}$ равен тождественному для инвариантного скалярного выхода и переставляет или поворачивает компоненты для эквивариантного поля. В разделе~\ref{sec:ml-map-data-loss} вход и выход задавались как пространства; здесь уточняется, какие геометрические записи являются элементами входного пространства и какие преобразования считаются служебными.

\medskip
\noindent\textbf{Параметрическое описание конструкции.}
В параметрическом представлении конструкция задаётся вектором
\begin{equation}
\label{eq:parametric-design-vector}
 \theta=(\theta_1,\ldots,\theta_m)\in\Theta\subseteq\R^m,
 \qquad
 \Omega=\Omega(\theta).
\end{equation}
Компонентами могут быть длины, толщины, радиусы, углы или коэффициенты сплайна. Отображение $\theta\mapsto\Omega(\theta)$ превращает изменение координат в изменение формы. Если оно дифференцируемо, правило цепочки даёт
\begin{equation}
 \dv{J}{\theta}
 =
 \pdv{J}{\Omega}
 \pdv{\Omega}{\theta},
\end{equation}
где первая производная обозначает чувствительность критерия к геометрии, а вторая --- чувствительность геометрии к параметрам. В вычислениях эти объекты реализуются через координаты узлов, неявную границу или отдельный геометрический оператор.

Преимущество параметризации состоит в фиксированной размерности и явных границах допустимости. Недостаток также следует из формулы~\eqref{eq:parametric-design-vector}: заранее выбранные координаты определяют доступное семейство форм. Если параметризация описывает только балку с одним отверстием, обучение или оптимизация не создадут второе отверстие. Изменение топологии требует другого пространства $\Theta$ либо другого представления.

Нормировка параметров выполняется до подачи в модель. Для границ $\theta_j^{\min}<\theta_j^{\max}$ используем
\begin{equation}
 \widetilde\theta_j
 =
 \frac{\theta_j-\theta_j^{\min}}
      {\theta_j^{\max}-\theta_j^{\min}}.
\end{equation}
Эта операция сохраняет физический параметр, но меняет координаты задачи обучения и её обусловленность, что связывает представление с разделом~\ref{sec:training-as-optimization}.

\subsection{Регулярная сетка и поле материала}
Для регулярной сетки конструкция записывается массивом значений
\begin{equation}
 \rho_h=(\rho_i)_{i=1}^{n_h},
 \qquad
 0\leq\rho_i\leq1,
\end{equation}
где индекс $i$ соответствует ячейке, а $h$ характеризует пространственное разрешение. Значение $\rho_i$ может описывать долю материала, фазу или другой локальный параметр модели. Такая запись естественно согласуется с тензорными операциями на регулярной решётке и позволяет изменять форму без ручного выбора геометрических размеров.

Однако массив $\rho_h$ является дискретизацией, а не самой конструкцией. При переходе от шага $h$ к $h/2$ меняются размер входа, число степеней свободы и минимальный разрешимый элемент формы. Поэтому равенство физических объектов нельзя определять покоординатным сравнением массивов разных размеров. Нужен оператор переноса
\begin{equation}
 \mathsf P_{h\to h'}:
 \mathcal X_h\to\mathcal X_{h'},
\end{equation}
после которого сравниваются геометрия, интегральные величины и физический отклик.

Пороговое преобразование $\rho_i\mapsto\mathbf 1_{\{\rho_i\geq\tau\}}$ создаёт бинарную конструкцию, но разрывает обычную дифференцируемость по $\rho_i$. Поэтому в обучении и проектной оптимизации различаются непрерывное поле, его дискретное представление и финальная геометрическая интерпретация. Градиент по непрерывным плотностям не следует автоматически считать градиентом по уже пороговой детали.

\medskip
\noindent\textbf{Множество точек и точечное облако.}
Точечное представление задаётся конечным множеством
\begin{equation}
 \mathcal P
 =
 \bigl\{(x_i,a_i)\bigr\}_{i=1}^{n},
 \qquad
 x_i\in\R^d,
 \quad
 a_i\in\R^q,
\end{equation}
где $x_i$ --- координата, а $a_i$ --- локальные признаки. Порядок перечисления точек не является частью геометрии. Если матрица перестановки $P$ меняет порядок строк $X$ и $A$, то скалярная модель обязана удовлетворять
\begin{equation}
 f_w(PX,PA)=f_w(X,A).
\end{equation}
Для выходного поля на тех же точках требуется эквивариантность
\begin{equation}
 F_w(PX,PA)=P F_w(X,A).
\end{equation}

Простейший инвариантный агрегатор имеет вид
\begin{equation}
 f_w(\mathcal P)
 =
 \psi_w\left(
 \sum_{i=1}^{n}
 \phi_w(x_i,a_i)
 \right).
\end{equation}
Сумма не зависит от порядка точек, но сама по себе не кодирует связность поверхности или объёма. Две выборки с близкими координатами могут описывать разные топологические объекты. Поэтому точечное облако удобно для измеренных или нерегулярных данных, но физическая постановка часто требует дополнительно восстановить соседство, нормали, границу и области приложения условий.

\subsection{Сетка как геометрический объект и как граф}
Для поверхностной или объёмной сетки используем расширенную запись
\begin{equation}
 \mathcal M=(V,E,F,X,H),
\end{equation}
где $V$, $E$ и $F$ задают комбинаторную структуру, $X\in\R^{|V|\times d}$ содержит координаты, а $H\in\R^{|V|\times q}$ --- признаки узлов. Для объёмной сетки к записи добавляются ячейки. Граф $G=(V,E)$ сохраняет соседство, но не обязательно хранит ориентацию граней, объём элементов и геометрические отображения, необходимые физическому решателю.

Локальное обновление узловых признаков можно записать как
\begin{equation}
 h_i^{(\ell+1)}
 =
 \Psi_{\ell}\left(
 h_i^{(\ell)},
 \sum_{j\in\mathcal N(i)}
 \Phi_{\ell}
 \bigl(h_i^{(\ell)},h_j^{(\ell)},x_j-x_i,e_{ij}\bigr)
 \right).
\end{equation}
Суммирование по соседям делает обновление независимым от порядка перечисления рёбер. Относительная координата $x_j-x_i$ устраняет зависимость от общего переноса, но ещё не гарантирует корректное поведение при повороте. Для этого необходимо отдельно определить, как преобразуются скалярные, векторные и тензорные признаки.

Перенумерация узлов задаётся матрицей перестановки $P$. Тогда
\begin{equation}
 X'=PX,
 \qquad
 H'=PH,
 \qquad
 A'=PAP^{\trans},
\end{equation}
где $A$ --- матрица смежности. Эквивариантная узловая модель удовлетворяет
\begin{equation}
 F_w(A',X',H')
 =
 P F_w(A,X,H).
\end{equation}
Это требование относится к нумерации. Сгущение или перестроение сетки меняет уже не только порядок, но и сам дискретный граф, поэтому перенос между сетками проверяется отдельно.

\subsection{Неявное поле и изменяемая топология}
Неявная геометрия задаётся скалярным полем
\begin{equation}
 \varphi_{\theta}:D\subseteq\R^d\to\R,
 \qquad
 \Omega(\theta)
 =
 \{x\in D:\varphi_{\theta}(x)<0\},
 \qquad
 \partial\Omega(\theta)
 =
 \{x:\varphi_{\theta}(x)=0\}.
\end{equation}
Параметры $\theta$ могут быть коэффициентами аналитической функции или весами обучаемого декодера. Одна формула описывает область без явного списка вершин, а изменение числа компонент или отверстий не требует заранее фиксировать комбинаторную структуру.

В регулярной точке границы, где $\nabla_x\varphi_{\theta}\neq0$, единичная нормаль равна
\begin{equation}
 n(x)
 =
 \frac{\nabla_x\varphi_{\theta}(x)}
      {\norm{\nabla_x\varphi_{\theta}(x)}_2}.
\end{equation}
Дифференцируемость поля не означает автоматической дифференцируемости всего расчёта. Для МКЭ неявную область требуется интегрировать или превратить в сетку; изменение пересечений, связности и числа элементов может создавать негладкие участки вычислительного графа. Поэтому отдельно проверяются производная поля, производная геометрического извлечения и производная физического критерия.

\subsection{Законченный расчёт: граница неявной окружности}
\label{subsec:worked-implicit-circle}
Рассмотрим семейство окружностей
\begin{equation}
 \varphi_r(x_1,x_2)
 =
 x_1^2+x_2^2-r^2,
 \qquad
 \Omega(r)=\{x:\varphi_r(x)<0\}.
\end{equation}
При $r=5$ точка $x=(3,4)$ лежит на границе, поскольку
\begin{equation}
 \varphi_5(3,4)
 =
 3^2+4^2-5^2
 =0.
\end{equation}
Градиент по пространственной координате и производная по параметру равны
\begin{equation}
 \nabla_x\varphi_5(3,4)
 =
 \begin{pmatrix}6\\8\end{pmatrix},
 \qquad
 \pdv{\varphi_r}{r}\bigg|_{r=5}
 =-10.
\end{equation}
Поэтому единичная внешняя нормаль в этой точке имеет вид
\begin{equation}
 n
 =
 \frac{1}{10}
 \begin{pmatrix}6\\8\end{pmatrix}
 =
 \begin{pmatrix}0.6\\0.8\end{pmatrix}.
\end{equation}

Проследим движение той же угловой точки границы:
\begin{equation}
 x(r)
 =
 r
 \begin{pmatrix}3/5\\4/5\end{pmatrix}.
\end{equation}
Тогда
\begin{equation}
 \dv{x}{r}
 =
 \begin{pmatrix}3/5\\4/5\end{pmatrix}
 =n.
\end{equation}
Увеличение радиуса на $\delta r=0.01$ перемещает границу в нормальном направлении приблизительно на $0.01$. Одновременно при фиксированной точке $(3,4)$ поле меняется на
\begin{equation}
 \varphi_{5.01}(3,4)-\varphi_5(3,4)
 =
 -2\cdot5\cdot0.01-(0.01)^2
 =-0.1001,
\end{equation}
а линейная оценка даёт $-10\cdot0.01=-0.1$. Расчёт показывает различие двух производных: $\partial_r\varphi$ измеряет изменение значения поля в фиксированной точке, а $\mathrm d x/\mathrm d r$ --- движение самой границы. В генеративном проектировании эти величины нельзя подменять друг другом.

\subsection{Симметрии, инвариантность и эквивариантность}
Пусть группа $G$ действует на входном пространстве оператором $\rho_{\mathcal X}(g)$ и на выходном пространстве оператором $\rho_{\mathcal Y}(g)$. Эквивариантность модели означает
\begin{equation}
 f_w\bigl(\rho_{\mathcal X}(g)x\bigr)
 =
 \rho_{\mathcal Y}(g)f_w(x).
\end{equation}
Инвариантность получается при $\rho_{\mathcal Y}(g)=I$. Здесь важно преобразовывать всю физическую задачу. Поворот одной геометрии при неподвижных нагрузках и закреплениях создаёт другую задачу, поэтому её податливость не обязана совпадать с исходной.

Для согласованного жёсткого движения $x\mapsto Qx+b$ скалярный критерий может быть инвариантен,
\begin{equation}
\label{eq:scalar-rigid-invariance}
 J(g\cdot c)=J(c),
\end{equation}
а векторное поле перемещений должно поворачиваться,
\begin{equation}
\label{eq:vector-field-rigid-equivariance}
 u_{g\cdot c}(Qx+b)
 =
 Q u_c(x).
\end{equation}
Формулы~\eqref{eq:scalar-rigid-invariance} и~\eqref{eq:vector-field-rigid-equivariance} задают разные требования к выходу одной и той же модели. Нельзя объявить сеть ``инвариантной к поворотам'', не уточнив, является выход скаляром, вектором или тензором.

\begin{connection}[с дифференциальной геометрией]
Групповое действие и эквивариантность продолжают язык главы~\ref{chap:manifolds-differential-geometry}. Координаты являются локальной записью объекта, а физически значимые величины должны преобразовываться согласованно при смене координат. В машинном обучении это условие переносится на входы, скрытые признаки и выходы вычислительного графа.
\end{connection}

\medskip
\noindent\textbf{Выбор представления для проектного контура.}
Единственного лучшего представления нет. Параметрический вектор удобен при фиксированном семействе форм и малом числе переменных. Регулярное поле допускает локальное изменение материала, но привязывает разрешение к сетке. Точечное облако сохраняет нерегулярные измерения, однако не содержит связность автоматически. Сетка и граф дают соседство и естественный интерфейс с дискретным решателем, но требуют контроля перенумерации и перестроения сетки. Неявное поле допускает изменение топологии, зато добавляет этап извлечения границы и сопряжения с физическим расчётом.

Перед обучением фиксируется шестичастный контракт
\begin{equation}
 \boxed{
 \begin{gathered}
 \text{физический объект}
 \to
 \text{кодирование}
 \\
 \text{служебные симметрии}
 \to
 \text{допустимые изменения формы}
 \\
 \text{интерфейс с решателем}
 \to
 \text{проверка переноса}
 \end{gathered}
 }.
\end{equation}
Проверка переноса включает новую нумерацию, жёсткое движение, другую плотность точек или сетки и, когда это допустимо, изменение топологии. Ошибка модели измеряется после приведения выходов к одному физическому объекту, а не между несогласованными массивами.

\begin{connection}[с AIRI и генеративным проектированием]
В полном контуре
\begin{equation}
 z
 \xmapsto{\ G_\phi\ }
 x
 \xmapsto{\ \mathsf D\ }
 c
 \xmapsto{\ \mathcal S\ }
 u
 \longmapsto
 J
\end{equation}
генератор выдаёт представление $x$, декодер $\mathsf D$ превращает его в конструкцию $c$, а физический решатель $\mathcal S$ вычисляет состояние. Представление определяет, какие кандидаты вообще достижимы, какие ограничения можно проверить до решателя и через какие операции пройдёт градиент. Именно этот интерфейс связывает машинное обучение с задачами AIRI по генеративному проектированию. Конкретный способ обучения генератора остаётся содержанием главы~9; здесь фиксируется математически корректный носитель генерируемой конструкции.
\end{connection}

Итог секции состоит в разделении объекта и его записи. Представление должно сохранять физически значимую информацию, устранять зависимость от служебной нумерации, согласованно преобразовывать поля и иметь определённый интерфейс с физическим решателем. Только после этого геометрические данные становятся корректным входом обучаемой модели.

\section{Суррогаты физических полей и обучаемые операторы}
\label{sec:surrogates-neural-operators}
В главах~\ref{chap:partial-differential-equations} и~\ref{chap:continuum-mechanics} физическая модель задавала состояние конструкции через дифференциальное уравнение и его дискретизацию. Здесь решатель не заменяется словом ``нейросеть'': сначала фиксируется оператор, который требуется приблизить, затем его дискретный интерфейс, функция ошибки и область допустимого применения суррогата.

Пусть $a\in\mathcal A$ объединяет коэффициенты уравнения, геометрию, нагрузку и граничные условия, а $u\in\mathcal U$ обозначает физическое поле. Непрерывная задача определяет оператор решения
\begin{equation}
\label{eq:solution-operator}
 \mathcal S:\mathcal A\to\mathcal U,
 \qquad
 a\longmapsto u=\mathcal S(a).
\end{equation}
После дискретизации на сетке с параметром $h$ возникает отображение
\begin{equation}
 S_h:A_h\to U_h,
 \qquad
 a_h\longmapsto U_h=S_h(a_h).
\end{equation}
Обычная модель на фиксированном массиве приближает конкретное $S_h$. Обучаемый оператор строится как единое семейство дискретных реализаций одного отображения между пространствами функций. Это различие определяет, что именно проверяется при смене сетки.

\subsection{Три уровня суррогатной модели}
Для одной конструкции можно обучать три разных объекта. Скалярный суррогат предсказывает критерий
\begin{equation}
 \widehat J=f_w(a_h)\in\R.
\end{equation}
Он подходит для быстрого ранжирования кандидатов, но не восстанавливает перемещения, напряжения или локальные ограничения. Полевая модель возвращает значения на выбранной дискретизации,
\begin{equation}
 \widehat U_h=F_{w,h}(a_h)\in U_h.
\end{equation}
Её размер и параметры могут быть привязаны к конкретной сетке. Операторная модель стремится задать
\begin{equation}
 \mathcal G_w:\mathcal A\to\mathcal U,
 \qquad
 \widehat u=\mathcal G_w(a),
\end{equation}
а на каждой сетке вычисляется дискретная реализация $G_{w,h}$ с теми же обучаемыми параметрами $w$.

Эти уровни не образуют автоматическую иерархию качества. Скалярная модель может точнее предсказывать податливость, чем полевая модель, если задача ограничена одним критерием. Напротив, для контроля напряжений одного числа недостаточно. Контракт выбирается по проектному решению, которое должно быть принято после предсказания.

\medskip
\noindent\textbf{Вход оператора: не только коэффициент поля.}
В записи~\eqref{eq:solution-operator} символ $a$ обозначает весь набор данных, определяющих задачу. Для линейной упругости его удобно раскрыть как
\begin{equation}
 a
 =
 \bigl(
  \Omega,
  C,
  f,
  t_N,
  g_D,
  \Gamma_N,
  \Gamma_D
 \bigr),
\end{equation}
где $C$ --- тензор упругости, $f$ --- объёмная сила, $t_N$ --- поверхностная нагрузка, а $g_D$ --- заданное перемещение. Если модель получает только поле материала, но нагрузка меняется между примерами, она не аппроксимирует однозначное отображение.

После выбора представления из раздела~\ref{sec:geometry-mesh-representation} вход преобразуется в дискретные признаки
\begin{equation}
 a_h
 =
 \mathsf R_h(a)
 =
 \bigl(X_h,H_h,B_h,F_h\bigr),
\end{equation}
где $X_h$ содержит координаты, $H_h$ --- материальные и геометрические признаки, $B_h$ --- код граничных условий, а $F_h$ --- нагрузки. Оператор дискретизации $\mathsf R_h$ является частью постановки. Без него фраза ``модель работает на разных сетках'' не определяет, как одна физическая задача превращается в разные входы.

\subsection{Полевая ошибка и производные физического поля}
Для узлового вектора перемещений естественная относительная ошибка имеет вид
\begin{equation}
 e_u
 =
 \frac{\norm{\widehat U_h-U_h}_{M_h}}
      {\norm{U_h}_{M_h}},
 \qquad
 \norm{V}_{M_h}^2=V^{\trans}M_hV,
\end{equation}
где $M_h$ --- массовая матрица или согласованная квадратурная матрица. Обычная евклидова норма зависит от числа узлов и плотности сетки, поэтому она не всегда сравнима между дискретизациями.

В механике малая ошибка перемещений не гарантирует малую ошибку напряжений. Деформация и напряжение содержат пространственную производную,
\begin{equation}
 \widehat\varepsilon_h
 =
 \nabla^s\widehat u_h,
 \qquad
 \widehat\sigma_h
 =
 C:\widehat\varepsilon_h.
\end{equation}
Высокочастотная ошибка небольшой амплитуды может слабо влиять на $\norm{\widehat u_h-u_h}_{L^2}$ и сильно влиять на $\nabla\widehat u_h$. Поэтому дополнительно используется энергетическая ошибка
\begin{equation}
 e_E
 =
 \frac{
  \left(
   \int_{\Omega}
   (\widehat\varepsilon_h-\varepsilon_h)
   :C:
   (\widehat\varepsilon_h-\varepsilon_h)
   \dd x
  \right)^{1/2}
 }{
  \left(
   \int_{\Omega}
   \varepsilon_h:C:\varepsilon_h
   \dd x
  \right)^{1/2}
 }.
\end{equation}
Максимальное напряжение проверяется отдельно, поскольку интегральная норма может скрыть локальный пик около отверстия, контакта или закрепления.

\subsection{Физические невязки и согласование единиц}
Для статической упругости предсказанное поле должно удовлетворять равновесию, граничным условиям и определяющему соотношению. Введём невязки
\begin{align}
 r_{\mathrm{eq}}
 &=
 -\operatorname{div}\widehat\sigma_h-f,
 \\
 r_N
 &=
 \widehat\sigma_h n-t_N,
 \\
 r_D
 &=
 \widehat u_h-g_D,
 \\
 r_C
 &=
 \widehat\sigma_h-C:\nabla^s\widehat u_h.
\end{align}
Они имеют разные размерности: $r_{\mathrm{eq}}$ измеряется как сила на объём, $r_N$ --- как сила на площадь, $r_D$ --- как длина, а $r_C$ --- как напряжение. Складывать их квадраты без масштабирования означает неявно задавать веса через выбранную систему единиц.

Пусть $L_0$, $U_0$ и $\Sigma_0$ --- характерные длина, перемещение и напряжение. Безразмерная функция потерь может быть записана как
\begin{equation}
\begin{aligned}
 \mathcal L_{\mathrm{phys}}
 ={}&
 \alpha
 \frac{L_0^2}{\Sigma_0^2|\Omega|}
 \norm{r_{\mathrm{eq}}}_{L^2(\Omega)}^2
 +
 \beta
 \frac{\norm{r_N}_{L^2(\Gamma_N)}^2}
      {\Sigma_0^2|\Gamma_N|}
 \\
 &+
 \gamma
 \frac{\norm{r_D}_{L^2(\Gamma_D)}^2}
      {U_0^2|\Gamma_D|}
 +
 \delta
 \frac{\norm{r_C}_{L^2(\Omega)}^2}
      {\Sigma_0^2|\Omega|}.
\end{aligned}
\end{equation}
Коэффициенты $\alpha,\beta,\gamma,\delta$ теперь управляют содержательным балансом условий, а не компенсируют случайные единицы измерения.

Полная цель обучения соединяет данные и физику,
\begin{equation}
 \mathcal L(w)
 =
 \mathcal L_{\mathrm{data}}(w)
 +
 \mathcal L_{\mathrm{phys}}(w).
\end{equation}
Нулевая невязка не доказывает совпадение с экспериментом, если исходная физическая модель неверна. Маленькая ошибка по данным также не доказывает равновесие вне обучающей выборки. Эти проверки отвечают на разные вопросы.

\subsection{Операторный слой и независимость параметров от сетки}
Один операторный слой можно представить как локальное линейное преобразование плюс нелокальный интегральный оператор:
\begin{equation}
 v_{\ell+1}(x)
 =
 \sigma_{\ell}\left(
 W_{\ell}v_{\ell}(x)
 +
 \int_D
 \kappa_{\ell,w}(x,y,a(x),a(y))
 v_{\ell}(y)
 \dd y
 +b_{\ell}(x)
 \right).
\end{equation}
Здесь $v_0(x)=P_w(a(x),x)$ является поднятием входных признаков, а финальное поле получается проекцией $\widehat u(x)=Q_w(v_L(x))$. Ядро $\kappa_{\ell,w}$ передаёт влияние между удалёнными точками и играет роль бесконечномерного аналога матрицы весов.

На сетке интеграл заменяется квадратурой,
\begin{equation}
 v_{\ell+1,i}
 =
 \sigma_{\ell}\left(
 W_{\ell}v_{\ell,i}
 +
 \sum_{j=1}^{n_h}
 \omega_{j,h}
 \kappa_{\ell,w}(x_{i,h},x_{j,h},a_{i,h},a_{j,h})
 v_{\ell,j}
 +b_{\ell}(x_{i,h})
 \right).
\end{equation}
Параметры $w$ не зависят от числа узлов, но результат зависит от координат и квадратурных весов. Простая возможность подать массив другой длины ещё не означает приближение одного непрерывного оператора. Требуется согласованность дискретной суммы с интегралом при $h\to0$.

\medskip
\noindent\textbf{Дискретизационная инвариантность как проверяемое свойство.}
Пусть $\mathsf R_h:\mathcal A\to A_h$ дискретизирует вход, а $\mathsf I_h:U_h\to\mathcal U$ восстанавливает поле. Семейство $G_{w,h}$ согласовано с непрерывным оператором $\mathcal G_w$, если
\begin{equation}
 \mathsf I_h
 G_{w,h}
 \mathsf R_h(a)
 \longrightarrow
 \mathcal G_w(a)
 \quad
 \text{при }h\to0
\end{equation}
для входов из заявленной области. На практике неизвестен точный предел, поэтому сравниваются последовательные сетки:
\begin{equation}
 d_{h,h/2}(a)
 =
 \frac{
  \norm{
   \mathsf I_h\widehat U_h
   -
   \mathsf I_{h/2}\widehat U_{h/2}
  }_{\mathcal U}
 }{
  \norm{
   \mathsf I_{h/2}\widehat U_{h/2}
  }_{\mathcal U}
 }.
\end{equation}
Одновременно вычисляется ошибка относительно физического решателя на каждой сетке. Малый $d_{h,h/2}$ сам по себе недостаточен: модель может устойчиво сходиться к неверному полю.

Протокол переноса содержит одну и ту же физическую задачу на сетках $h$, $h/2$ и $h/4$, независимую перестройку сетки и хотя бы одну дискретизацию, не использованную при обучении. Проверяются поле, интегральный критерий и локальные экстремумы. Так понятие из статьи об обучаемых операторах превращается в воспроизводимый тест, а не в название архитектуры.

\subsection{Разложение полной ошибки}
Пусть $u$ --- решение непрерывной физической задачи, $u_h=\mathsf I_hS_h\mathsf R_h(a)$ --- результат численного решателя, а $\widehat u_h=\mathsf I_hG_{w,h}\mathsf R_h(a)$ --- предсказание. Неравенство треугольника даёт
\begin{equation}
 \norm{\widehat u_h-u}_{\mathcal U}
 \leq
 \underbrace{
  \norm{\widehat u_h-u_h}_{\mathcal U}
 }_{\text{ошибка суррогата относительно решателя}}
 +
 \underbrace{
  \norm{u_h-u}_{\mathcal U}
 }_{\text{ошибка дискретизации физической модели}}.
\end{equation}
Если данные получены решателем, обучение контролирует первое слагаемое, но не устраняет второе. При сравнении с экспериментом добавляется ошибка самой физической модели и измерений. Поэтому формулировка ``суррогат точнее решателя'' требует отдельного эталона, а не сравнения с теми же синтетическими метками.

Для проектного критерия $J(u,a)$ при локальной липшицевости
\begin{equation}
 \abs{J(\widehat u_h,a)-J(u,a)}
 \leq
 L_J
 \norm{\widehat u_h-u}_{\mathcal U}.
\end{equation}
Константа $L_J$ зависит от критерия. Для максимального напряжения она может быть значительно хуже, чем для интегральной энергии. Поэтому полевая метрика выбирается с учётом последующего проектного решения.

\subsection{Законченный расчёт: точное значение и неверное поле}
\label{subsec:worked-bar-surrogate}
Рассмотрим однородный стержень длины $L$, площади $A$ и модуля Юнга $E$. Левый конец закреплён, а к правому приложена сила $F$. Точное решение задачи
\begin{equation}
 -\dv{}{x}\left(EA\dv{u}{x}\right)=0,
 \qquad
 u(0)=0,
 \qquad
 EA\dv{u}{x}(L)=F
\end{equation}
равно
\begin{equation}
 u(x)
 =
 u_L\frac{x}{L},
 \qquad
 u_L=\frac{FL}{EA}.
\end{equation}
Пусть суррогат выдал поле
\begin{equation}
 \widehat u(x)
 =
 u_L\left[
 \frac{x}{L}
 +
 c\frac{x}{L}\left(1-\frac{x}{L}\right)
 \right].
\end{equation}
Оно точно воспроизводит оба перемещения на концах:
\begin{equation}
 \widehat u(0)=0,
 \qquad
 \widehat u(L)=u_L.
\end{equation}
Следовательно, скалярная ошибка по перемещению нагруженного конца равна нулю.

Относительная $L^2$-ошибка поля вычисляется заменой $s=x/L$:
\begin{equation}
 \frac{\norm{\widehat u-u}_{L^2(0,L)}}
      {\norm{u}_{L^2(0,L)}}
 =
 |c|
 \sqrt{
  \frac{\int_0^1s^2(1-s)^2\dd s}
       {\int_0^1s^2\dd s}
 }
 =
 \frac{|c|}{\sqrt{10}}.
\end{equation}
Для $c=0.2$ это $0.0632$, то есть $6.32\%$. Энергетическая ошибка равна
\begin{equation}
 \frac{
  \norm{\widehat u'-u'}_{L^2(0,L)}
 }{
  \norm{u'}_{L^2(0,L)}
 }
 =
 |c|
 \sqrt{
  \int_0^1(1-2s)^2\dd s
 }
 =
 \frac{|c|}{\sqrt3}
 =0.1155.
\end{equation}

Предсказанное усилие имеет вид
\begin{equation}
 \widehat N(x)
 =
 EA\widehat u'(x)
 =
 F\left[
 1+c\left(1-2\frac{x}{L}\right)
 \right].
\end{equation}
Поэтому внутренняя невязка равновесия и ошибка правой граничной силы равны
\begin{equation}
 -\widehat N'(x)
 =
 \frac{2cF}{L},
 \qquad
 \widehat N(L)-F
 =
 -cF.
\end{equation}
При $F=100\,\mathrm N$, $L=2\,\mathrm m$ и $c=0.2$ получаем постоянную невязку $20\,\mathrm{N/m}$ и граничную ошибку $-20\,\mathrm N$. Безразмерные квадратичные штрафы равны
\begin{equation}
 \frac{L^2}{F^2L}
 \int_0^L
 \left(-\widehat N'(x)\right)^2\dd x
 =4c^2=0.16,
 \qquad
 \left(\frac{\widehat N(L)-F}{F}\right)^2
 =c^2=0.04.
\end{equation}
Расчёт показывает, почему точное предсказание одного проектного числа не сертифицирует поле. Суррогат правильно воспроизвёл перемещение конца, но нарушил равновесие и приложенную нагрузку.

\medskip
\noindent\textbf{Протокол применения в генеративном проектировании.}
Суррогат может использоваться для предварительного отбора большого числа кандидатов, для приближённого вычисления поля или как дифференцируемая часть проектного контура. В каждом случае сохраняется физический решатель, но меняется частота его вызова. Быстрый предиктор не становится источником окончательной сертификации конструкции.

Рабочий контур имеет вид
\begin{equation}
 \theta
 \xmapsto{\ \mathsf E\ }
 a_h
 \xmapsto{\ G_{w,h}\ }
 \widehat U_h
 \longmapsto
 \widehat J,
 \qquad
 \text{отобранные кандидаты}
 \xmapsto{\ S_h\ }
 U_h.
\end{equation}
Для отобранных решений повторно проверяются равновесие, граничные условия, напряжения и критерии на независимой сетке. Если суррогат используется для градиента по $\theta$, сравнивается не только $\widehat J$, но и направление производной
\begin{equation}
 \cos\alpha
 =
 \frac{
  \left\langle
   \nabla_{\theta}\widehat J,
   \nabla_{\theta}J_h
  \right\rangle
 }{
  \norm{\nabla_{\theta}\widehat J}_2
  \norm{\nabla_{\theta}J_h}_2
 }.
\end{equation}
Малая ошибка значения при отрицательном $\cos\alpha$ способна направить оптимизацию в противоположную сторону.

\begin{connection}[с AIRI и генеративным проектированием]
Для задач AIRI суррогатный оператор полезен не как абстрактная замена механики, а как ускоряемый интерфейс между представлением конструкции и физическим откликом. Он позволяет оценивать много кандидатов, переносить поля между дискретизациями и получать производные внутри обучаемого контура. При этом физический решатель остаётся источником данных, независимой проверки и финального решения о допустимости. Глава~9 будет использовать этот интерфейс при построении генератора; текущий раздел фиксирует, какие выходы суррогата можно передавать в проектный цикл и какие проверки нельзя исключать.
\end{connection}

Итог секции состоит в разделении трёх объектов: непрерывного оператора $\mathcal S$, численного решателя $S_h$ и обучаемого приближения $G_{w,h}$. Качество определяется не одним средним числом, а ошибкой поля, производных величин, физических невязок, переносом между сетками и корректностью проектного критерия.

\section{Гибридная физико-обучаемая модель и протокол проверки}
\label{sec:hybrid-physics-ml-validation}
Гибридной будем называть модель, в которой обучаемое отображение и физический оператор входят в один вычислительный контур, но сохраняют разные роли. Обучаемый блок может восстанавливать неизвестную часть закона материала, приближать дорогое решение, оценивать скрытые параметры по наблюдениям или порождать описание конструкции. Физический блок задаёт состояние, ограничения и критерии, которые следуют из выбранной модели среды. Смешивать эти роли нельзя: ошибка обучения, ошибка дискретизации и несоответствие физической модели имеют разные причины и проверяются разными способами.

Сквозной контур главы записывается как
\begin{equation}
 z
 \xmapsto{\ G_\phi\ }
 \theta
 \xmapsto{\ \mathsf E\ }
 a_h
 \xmapsto{\ \mathcal H_{w,h}\ }
 \widehat U_h
 \longmapsto
 \widehat J_h,
\end{equation}
где $z$ содержит условие или латентный код, $G_\phi$ формирует параметры конструкции, $\mathsf E$ кодирует её для вычислений, а $\mathcal H_{w,h}$ обозначает гибридный физико-обучаемый оператор. В частном случае $\mathcal H_{w,h}=S_h$ и обучаемый блок отсутствует внутри решателя. В другом частном случае $\mathcal H_{w,h}=G_{w,h}$ является суррогатом из раздела~\ref{sec:surrogates-neural-operators}. Наиболее содержательный случай возникает, когда нейросеть определяет только один компонент физической модели, а состояние по-прежнему вычисляется из уравнений равновесия.

\subsection{Четыре роли обучаемого блока}
Первую роль задаёт прямой суррогат
\begin{equation}
 a_h\longmapsto \widehat U_h
 \quad\text{или}\quad
 a_h\longmapsto \widehat J_h.
\end{equation}
Он ускоряет многократную оценку кандидатов, но не изменяет саму физическую постановку. Его область применимости определяется данными и проверкой из раздела~\ref{sec:surrogates-neural-operators}.

Вторая роль состоит в идентификации неизвестных параметров по наблюдениям:
\begin{equation}
 y_{\mathrm{obs}}
 \xmapsto{\ q_w\ }
 \widehat\mu,
 \qquad
 S_h(\widehat\mu)\approx y_{\mathrm{obs}}.
\end{equation}
Здесь выход сети не принимается сам по себе. Он подставляется в прямую физическую модель, после чего сравнивается восстановленный отклик с наблюдениями.

Третья роль --- обучаемый компонент внутри решателя. Например, вместо заранее заданного конститутивного закона используется
\begin{equation}
 \sigma
 =
 \mathcal C_w(\varepsilon,\xi),
\end{equation}
где $\xi$ содержит состояние материала или дополнительные признаки. Баланс импульса и граничные условия при этом не исчезают. Они определяют, какие значения $\varepsilon$ и $\sigma$ допустимы в конструкции.

Четвёртая роль относится к генератору
\begin{equation}
 G_\phi: \mathcal Z\times\mathcal C\to\Theta,
\end{equation}
который по коду $z$ и условию $c$ выдаёт параметры конструкции. Эта роль будет развита в главе~9. В текущей главе фиксируется интерфейс генератора с физикой: его выход должен быть кодируемым, решаемым и проверяемым.

\medskip
\noindent\textbf{Функция обучения и проектный критерий.}
Параметры модели $w$ определяются по функции обучения
\begin{equation}
 w^\star
 =
 \argmin_{w\in\mathcal W}
 \mathcal L_{\mathrm{train}}(w).
\end{equation}
Параметры конструкции $\theta$ определяются по другому объекту:
\begin{equation}
 \theta^\star
 =
 \argmin_{\theta\in\Theta_{\mathrm{adm}}}
 J_h\bigl(\theta,U_h(\theta)\bigr).
\end{equation}
Даже если $\mathcal L_{\mathrm{train}}$ и $J_h$ используют одну физическую величину, они не совпадают. Первая функция измеряет качество модели на наборе примеров; вторая сравнивает конструкции и включает проектные ограничения.

В гибридном контуре суррогатная цель имеет вид
\begin{equation}
 \widehat J_h(\theta)
 =
 J_h\bigl(\theta,\mathcal H_{w,h}(\mathsf E(\theta))\bigr).
\end{equation}
Малое значение $\mathcal L_{\mathrm{train}}(w)$ не гарантирует, что минимум $\widehat J_h$ расположен рядом с минимумом $J_h$. Оптимизация специально ищет области, где критерий улучшается, и может вывести входы за пределы плотной части обучающей выборки. Поэтому проверка суррогата на случайных тестовых примерах дополняется проверкой на траектории проектирования.

\subsection{Дифференцирование полного гибридного графа}
Пусть дискретное состояние определяется неявным уравнением
\begin{equation}
 R_h(U,\theta,w)=0.
\end{equation}
Если матрица $\partial_U R_h$ обратима, дифференцирование по параметру конструкции даёт
\begin{equation}
 \partial_U R_h\,\frac{\dd U}{\dd\theta}
 +
 \partial_\theta R_h
 +
 \partial_w R_h\,\frac{\dd w}{\dd\theta}
 =0.
\end{equation}
При фиксированных обученных весах $w$ последнее слагаемое равно нулю. Однако зависимость обучаемого компонента от $\theta$ через его вход уже содержится в $\partial_\theta R_h$. Удаление этой ветви вычислительного графа даёт неполную производную.

Для скалярного критерия $J_h(\theta,U)$ вводится сопряжённый вектор $\lambda$, определённый системой
\begin{equation}
 \bigl(\partial_U R_h\bigr)^{\trans}\lambda
 =
 \bigl(\partial_U J_h\bigr)^{\trans}.
\end{equation}
Тогда полная производная равна
\begin{equation}
 \frac{\dd J_h}{\dd\theta}
 =
 \partial_\theta J_h
 -
 \lambda^{\trans}\partial_\theta R_h.
\end{equation}
Эта формула объединяет правило цепочки из главы~\ref{chap:multivariable-calculus}, обратный режим из раздела~\ref{sec:computational-graph-autodiff} и сопряжённый подход физических глав. Автоматическое дифференцирование вычисляет локальные произведения с якобианами, но не решает за исследователя, какие зависимости должны присутствовать в $R_h$.

\subsection{Обучаемый закон внутри физической модели}
Если сеть задаёт конститутивное отображение, её выход должен удовлетворять требованиям выбранной механической модели. Для гиперупругого материала безопаснее параметризовать плотность энергии
\begin{equation}
 \Psi_w(F,\xi),
 \qquad
 P_w
 =
 \frac{\partial\Psi_w}{\partial F},
\end{equation}
чем независимо предсказывать компоненты напряжения. Такая запись автоматически связывает напряжение с одной потенциальной функцией. Она, однако, сама по себе не гарантирует объективность, положительность энергии или устойчивость. Эти свойства должны быть встроены в представление либо проверены на допустимом диапазоне деформаций.

Для малых деформаций обучаемая поправка может быть отделена от базового закона:
\begin{equation}
 \sigma_w(\varepsilon)
 =
 C_0:\varepsilon
 +
 \Delta\sigma_w(\varepsilon).
\end{equation}
Базовая часть фиксирует известный предел, а сеть аппроксимирует только необъяснённую составляющую. Такая декомпозиция уменьшает объём задачи обучения, но не превращает поправку в физически допустимую автоматически. Проверяются симметрия касательного оператора, диапазон деформаций и отсутствие отрицательной локальной жёсткости там, где модель материала её не допускает.

\subsection{Разделение данных без утечки физического сценария}
В обычной выборке соседние строки могут соответствовать почти одной и той же конструкции, сетке или нагрузке. Случайное разделение таких строк создаёт утечку: тестовый пример формально новый, но физически почти совпадает с обучающим. Поэтому единицей разделения служит не отдельный вектор, а сценарий
\begin{equation}
 s
 =
 \bigl(
  \text{семейство геометрии},
  \text{материал},
  \text{нагрузка},
  \text{граничные условия}
 \bigr).
\end{equation}
Все дискретизации и близкие вариации одного сценария помещаются в одну часть данных.

Для задачи генеративного проектирования требуется ещё более строгая проверка. Тестовый набор должен содержать не только случайные конструкции, но и кандидаты, полученные самим алгоритмом оптимизации или генератором. Они принадлежат распределению применения, которое может заметно отличаться от распределения исходной выборки. Это различие уже было введено в разделе~\ref{sec:ml-map-data-loss}; здесь оно становится частью рабочего протокола.

\subsection{Четыре уровня проверки}
Проверка по данным сравнивает $\widehat y$ с эталонным $y$ на независимых сценариях. Она отвечает на вопрос, воспроизводит ли модель имеющиеся расчёты или измерения.

Физическая проверка оценивает невязки уравнений, граничных условий, симметрий и материальных ограничений. Она отвечает на вопрос, совместимо ли предсказание с выбранной моделью среды.

Численная проверка сравнивает результат при сгущении сетки, изменении шага интегрирования и другой допустимой дискретизации. Она отделяет поведение обучаемого блока от артефактов конкретной реализации $S_h$.

Проектная проверка заново вычисляет критерии и ограничения для кандидата независимым решателем. Она отвечает на конечный вопрос: можно ли принять конструкцию, а не только её предсказанное поле. Эти четыре уровня нельзя свести к одной усреднённой метрике.

Для кандидата $\theta$ удобно хранить вектор проверки
\begin{equation}
 \mathcal V(\theta)
 =
 \bigl(
  e_{\mathrm{data}},
  r_{\mathrm{phys}},
  e_{\mathrm{mesh}},
  g_{\mathrm{design}}
 \bigr),
\end{equation}
где $g_{\mathrm{design}}$ содержит значения проектных ограничений. Порог каждого компонента задаётся до запуска массового поиска, а не подбирается после просмотра лучших результатов.

\medskip
\noindent\textbf{Индикатор выхода из области применимости.}
Ни одна скалярная величина не доказывает, что вход находится внутри области обучения. Практический индикатор можно собрать из трёх независимых сигналов:
\begin{equation}
 \eta(\theta)
 =
 \alpha_d d_{\mathrm{train}}(\theta)
 +
 \alpha_r r_{\mathrm{phys}}(\theta)
 +
 \alpha_e s_{\mathrm{ens}}(\theta).
\end{equation}
Здесь $d_{\mathrm{train}}$ измеряет расстояние до известных сценариев в выбранном представлении, $r_{\mathrm{phys}}$ --- физическую невязку, а $s_{\mathrm{ens}}$ --- разброс нескольких независимо обученных моделей. Коэффициенты $\alpha_d,\alpha_r,\alpha_e$ только приводят компоненты к сопоставимому масштабу.

Индикатор $\eta$ не является вероятностью отказа и не заменяет физический расчёт. Он выполняет роль переключателя:
\begin{equation}
 \begin{aligned}
  \eta(\theta)\leq\eta_{\max}
  &\quad\Longrightarrow\quad
  \text{разрешена быстрая оценка},\\
  \eta(\theta)>\eta_{\max}
  &\quad\Longrightarrow\quad
  \text{обязателен вызов }S_h.
 \end{aligned}
\end{equation}
Новый точный расчёт добавляется в набор данных только после проверки качества сетки и постановки. Так возникает управляемый цикл расширения данных, а не бесконтрольное самообучение на собственных предсказаниях.

\subsection{Законченный расчёт: потерянная ветвь гибридного градиента}
\label{subsec:worked-hybrid-gradient}
Рассмотрим систему с двумя степенями свободы
\begin{equation}
 K(\theta)U=F,
 \qquad
 F=\begin{pmatrix}1\\0\end{pmatrix},
\end{equation}
где обучаемая поправка увеличивает первый диагональный элемент:
\begin{equation}
 K(\theta)
 =
 \begin{pmatrix}
  1+\theta+q_w(\theta) & -1\\
  -1 & 2
 \end{pmatrix},
 \qquad
 q_w(\theta)=0.2\theta.
\end{equation}
При $\theta=1$ имеем
\begin{equation}
 K(1)
 =
 \begin{pmatrix}
  2.2 & -1\\
  -1 & 2
 \end{pmatrix},
 \qquad
 U(1)
 =
 \begin{pmatrix}
  10/17\\5/17
 \end{pmatrix}.
\end{equation}
Пусть критерий равен податливости
\begin{equation}
 J(\theta)=F^{\trans}U(\theta)=U_1(\theta).
\end{equation}
Поскольку
\begin{equation}
 J(\theta)
 =
 \frac{2}{1+2.4\theta},
\end{equation}
точная производная при $\theta=1$ равна
\begin{equation}
 J'(1)
 =
 -\frac{4.8}{(3.4)^2}
 =
 -\frac{120}{289}
 \approx
 -0.4152.
\end{equation}
Та же величина получается из формулы чувствительности
\begin{equation}
 J'(1)
 =
 -U(1)^{\trans}
 \frac{\partial K}{\partial\theta}(1)
 U(1),
 \qquad
 \frac{\partial K}{\partial\theta}(1)
 =
 \begin{pmatrix}
  1.2 & 0\\0 & 0
 \end{pmatrix}.
\end{equation}
Если ошибочно считать выход $q_w(\theta)$ константой при обратном проходе, используется матрица
\begin{equation}
 \left.\frac{\partial K}{\partial\theta}\right|_{q_w\ \mathrm{detached}}
 =
 \begin{pmatrix}
  1 & 0\\0 & 0
 \end{pmatrix}
\end{equation}
и получается
\begin{equation}
 \widetilde J'(1)
 =
 -\frac{100}{289}
 \approx
 -0.3460.
\end{equation}
Относительная ошибка модуля градиента составляет
\begin{equation}
 \frac{
  \abs{\widetilde J'(1)-J'(1)}
 }{
  \abs{J'(1)}
 }
 =
 \frac16
 \approx
 16.7\%.
\end{equation}
Значение состояния было вычислено из правильной матрицы, но производная оказалась неверной из-за одной удалённой ветви графа. Поэтому проверка гибридной модели включает конечную разность полного контура, а не только тесты отдельных нейросетевых операций.

\subsection{Рабочий контур для AIRI и генеративного проектирования}
Полный цикл начинается с условия задачи $c$: области допустимых размеров, материала, нагрузок и ограничений. Генератор или оптимизатор предлагает $\theta$. Кодировщик $\mathsf E$ проверяет представимость и строит вход решателя. Быстрый гибридный оператор возвращает поле, критерии, физическую невязку и индикатор $\eta$. Только кандидаты, прошедшие дешёвые проверки, направляются в независимый физический расчёт. Результат этого расчёта используется одновременно для принятия конструкции и для пополнения данных в недостаточно покрытой области.

Этот контур можно записать как
\begin{equation}
 c,z
 \longmapsto
 \theta
 \longmapsto
 \bigl(
  \widehat U_h,
  \widehat J_h,
  r_{\mathrm{phys}},
  \eta
 \bigr)
 \longmapsto
 \text{отбор}
 \xmapsto{\ S_{h/2}\ }
 \bigl(U_{h/2},J_{h/2},g_{h/2}\bigr).
\end{equation}
Именно такая запись связывает главу с задачами AIRI без неподтверждённых утверждений о конкретных внутренних системах: машинное обучение ускоряет представление, предсказание и поиск, а математическая физика задаёт проверяемые состояния и ограничения.

\begin{connection}[с главой о генеративном проектировании]
Глава~9 начнётся не с повторного объяснения нейросети, а с уже готового интерфейса. Из текущей главы она получает пространство параметров $\Theta$, кодировщик $\mathsf E$, обучаемый или гибридный оператор $\mathcal H_{w,h}$, проектный критерий $J_h$, градиент и протокол независимой проверки. Новым объектом станет распределение или механизм порождения кандидатов. Поэтому генеративность будет означать способ исследовать множество конструкций, а не отмену физической модели.
\end{connection}

Итог главы можно выразить одной цепочкой:
\begin{equation}
 \begin{gathered}
  \text{математический контракт}
  \to \text{данные}
  \to \text{обучение},\\
  \text{геометрическое представление}
  \to \text{физический или гибридный оператор},\\
  \text{независимая проверка}.
 \end{gathered}
\end{equation}
Машинное обучение становится частью проектирования только тогда, когда каждый переход в этой цепочке определён и проверяем. Точность по данным, выполнение физических законов и пригодность конструкции остаются разными утверждениями.
